% =============================================================
%  Documento di studio per l'orale di Elementi di Intelligenza
%  Artificiale (6 CFU) — Progetto: NAS via Algoritmo Genetico
%  Autore: Giovanni
% =============================================================
\documentclass[11pt,a4paper]{article}

% ─── Lingua, codifica, font ──────────────────────────────────
\usepackage[utf8]{inputenc}
\usepackage[T1]{fontenc}
\usepackage[main=italian,provide=*]{babel}

% ─── Layout ──────────────────────────────────────────────────
\usepackage[margin=2.5cm]{geometry}
\usepackage{setspace}
\setstretch{1.15}
\usepackage{parskip}
\usepackage{titlesec}
\titleformat{\section}{\Large\bfseries\color{darkblue}}{\thesection}{1em}{}
\titleformat{\subsection}{\large\bfseries\color{midblue}}{\thesubsection}{1em}{}
\titleformat{\subsubsection}{\normalsize\bfseries}{\thesubsubsection}{1em}{}

% ─── Matematica ─────────────────────────────────────────────
\usepackage{amsmath, amssymb, amsthm}
\usepackage{mathtools}
\usepackage{bm}

% ─── Tabelle ─────────────────────────────────────────────────
\usepackage{booktabs}
\usepackage{array}
\usepackage{colortbl}
\usepackage{multirow}

% ─── Colori ──────────────────────────────────────────────────
\usepackage[table,dvipsnames]{xcolor}
\definecolor{darkblue}{HTML}{1B3A57}
\definecolor{midblue}{HTML}{2A6F97}
\definecolor{headerblue}{HTML}{2A6F97}
\definecolor{rowgray}{HTML}{F2F4F7}
\definecolor{accent}{HTML}{C1666B}
\definecolor{boxbg}{HTML}{F7F9FC}
\definecolor{conceptcolor}{HTML}{4F8C6C}
\definecolor{questioncolor}{HTML}{8C4F4F}
\definecolor{warncolor}{HTML}{B5651D}

% ─── Pseudocodice ────────────────────────────────────────────
\usepackage{algorithm}
\usepackage{algpseudocode}
\renewcommand{\algorithmicrequire}{\textbf{Input:}}
\renewcommand{\algorithmicensure}{\textbf{Output:}}
\floatname{algorithm}{Algoritmo}

% ─── Box decorativi ──────────────────────────────────────────
\usepackage[most]{tcolorbox}

% Box "concetto chiave"
\newtcolorbox{concetto}[1][]{
    enhanced,
    breakable,
    colback=conceptcolor!5,
    colframe=conceptcolor!70,
    fonttitle=\bfseries,
    title={Concetto chiave: #1},
    boxrule=0.6pt,
    arc=2pt,
    left=8pt, right=8pt, top=6pt, bottom=6pt,
}

% Box "domanda d'esame"
\newtcolorbox{domanda}{
    enhanced,
    breakable,
    colback=questioncolor!5,
    colframe=questioncolor!70,
    fonttitle=\bfseries,
    title={Domanda tipo orale},
    boxrule=0.6pt,
    arc=2pt,
    left=8pt, right=8pt, top=6pt, bottom=6pt,
}

% Box "attenzione/insidia"
\newtcolorbox{attenzione}{
    enhanced,
    breakable,
    colback=warncolor!5,
    colframe=warncolor!70,
    fonttitle=\bfseries,
    title={Attenzione},
    boxrule=0.6pt,
    arc=2pt,
    left=8pt, right=8pt, top=6pt, bottom=6pt,
}

% Box "riassunto sezione"
\newtcolorbox{riassunto}{
    enhanced,
    breakable,
    colback=midblue!5,
    colframe=midblue!50,
    fonttitle=\bfseries,
    title={Punti chiave da ricordare},
    boxrule=0.6pt,
    arc=2pt,
    left=8pt, right=8pt, top=6pt, bottom=6pt,
}

% ─── Hyperref (sempre per ultimo nei pacchetti) ──────────────
\usepackage[
    colorlinks=true,
    linkcolor=darkblue,
    citecolor=darkblue,
    urlcolor=midblue,
]{hyperref}

% ─── Comandi personalizzati ──────────────────────────────────
\newcommand{\fnt}{\texttt}     % font per nomi di funzioni/parametri
\newcommand{\param}[1]{\texttt{#1}}

% =============================================================
\begin{document}

% ─── FRONTESPIZIO ────────────────────────────────────────────
\begin{titlepage}
    \centering
    \vspace*{2cm}
    {\Huge\bfseries\color{darkblue} Neural Architecture Search\\[0.3cm]
    via Algoritmo Genetico\par}
    \vspace{0.5cm}
    {\large\itshape Documento di studio per l'orale\par}
    \vspace{2cm}
    {\Large Esame: Elementi di Intelligenza Artificiale (6 CFU)\par}
    \vspace{0.3cm}
    {\large Corso di laurea in Ingegneria Informatica\par}
    \vfill
    \begin{tcolorbox}[colback=boxbg, colframe=midblue!50, boxrule=0.5pt]
    \textbf{Cosa fa il progetto in una frase}\\[0.2cm]
    Un Algoritmo Genetico cerca automaticamente la migliore architettura di una rete
    neurale (numero di layer, neuroni per layer, funzione di attivazione) per un dato
    dataset, valutando ogni architettura tramite l'accuratezza ottenuta dopo training
    con backpropagation, penalizzata per la complessità della rete.
    \end{tcolorbox}
    \vfill
    {\large \today\par}
\end{titlepage}

% ─── INDICE ──────────────────────────────────────────────────
\tableofcontents
\newpage

% =============================================================
\section{Introduzione}
% =============================================================

% ─────────────────────────────────────────────────────────────
\subsection{L'idea del progetto in 60 secondi}
% ─────────────────────────────────────────────────────────────

Quando si progetta una rete neurale, la scelta dell'architettura
(\emph{quanti layer? quanti neuroni? quale funzione di attivazione?}) è
quasi sempre fatta a mano, per tentativi ed errori, basandosi sull'esperienza
del progettista. Il problema è che lo spazio delle architetture possibili è
enorme: anche limitandosi a $1\dots 5$ layer, $4\dots 64$ neuroni per layer e
$5$ funzioni di attivazione, si ottengono milioni di combinazioni. Provarle
tutte è impensabile.

L'idea del progetto è semplice e potente: invece di cercare a mano,
\textbf{lasciamo cercare a un Algoritmo Genetico (GA)}. Il GA imita
l'evoluzione naturale: parte da una popolazione casuale di architetture,
seleziona le migliori, le combina (crossover), introduce piccole variazioni
(mutazione) e ripete il processo per molte generazioni. Quello che emerge
alla fine è un'architettura quasi ottima, scoperta automaticamente.

I pesi della rete vengono comunque imparati con \textbf{backpropagation
classica}: il GA decide solo la \emph{forma} della rete, la backpropagation
ne riempie i \emph{valori}. Si tratta quindi di un'\textbf{ottimizzazione
annidata}: due ottimizzatori operano contemporaneamente su scale diverse.

% ─────────────────────────────────────────────────────────────
\subsection{Cos'è il Neural Architecture Search (NAS)}
% ─────────────────────────────────────────────────────────────

Il \textbf{Neural Architecture Search} è il sottocampo dell'\emph{automated
machine learning} che studia come automatizzare la progettazione delle reti
neurali. È strettamente collegato a:
\begin{itemize}
    \item \textbf{Ottimizzazione degli iperparametri} — gli iperparametri
          sono tutti i parametri di un modello che \emph{non} vengono
          appresi dai dati (learning rate, topologia, dimensione).
          Tradizionalmente si scelgono con \emph{grid search}, \emph{random
          search} o intuizione. NAS lo fa con tecniche evolutive o di
          reinforcement learning.
    \item \textbf{Meta-apprendimento} — ``imparare ad apprendere''. Un vero
          sistema di meta-learning estrae \emph{principi generali} validi
          su molti dataset (es. ``problemi con poche classi e feature
          continue beneficiano di pochi layer ReLU''). Questo progetto è
          più correttamente \emph{architecture search} per un singolo
          dataset, non meta-learning in senso stretto.
\end{itemize}

NAS è un campo industrialmente rilevante: modelli come \emph{EfficientNet}
(Google) sono nati da ricerca automatica dell'architettura, e NAS è la
tecnica di riferimento quando bisogna deployare reti su hardware con
risorse limitate (mobile, IoT).

% ─────────────────────────────────────────────────────────────
\subsection{Originalità e rilevanza}
% ─────────────────────────────────────────────────────────────

\begin{domanda}
\textit{``Il tuo progetto è originale?''}\\
No, nel senso accademico stretto: la NAS via GA è una tecnica nota dagli
anni '90. Ciò che è originale è la \emph{combinazione specifica} qui
realizzata:
\begin{itemize}
    \item rete neurale scritta da zero con backpropagation manuale (solo
          \fnt{numpy}, niente TensorFlow/PyTorch);
    \item GA scritto da zero che ottimizza l'architettura;
    \item penalità di complessità logaritmica per evitare reti grandi inutili;
    \item valutazione multipla ($K=3$) per ridurre la varianza dovuta
          all'inizializzazione casuale dei pesi.
\end{itemize}
Tutti i componenti sono costruiti senza librerie di machine learning
preconfezionate: questo dimostra la comprensione \emph{sotto il cofano}
dei due paradigmi --- reti neurali e algoritmi evolutivi --- più che la
capacità di usare API.
\end{domanda}

\begin{riassunto}
\begin{itemize}
    \item Il progetto risolve il problema della scelta manuale dell'architettura di una rete neurale.
    \item Usa un GA per esplorare lo spazio delle architetture e la backpropagation per addestrare ogni rete candidata.
    \item È un caso di \emph{ottimizzazione annidata}: due livelli di ottimizzazione, GA sopra e backpropagation sotto.
    \item Il valore non è nell'originalità del metodo, ma nell'implementazione completa da zero di entrambi i paradigmi.
\end{itemize}
\end{riassunto}

\newpage

% =============================================================
\section{Background di Machine Learning}
% =============================================================

Per inquadrare il progetto serve un nucleo solido di concetti di machine
learning. Tutto ciò che segue verrà richiamato nel resto del documento.

% ─────────────────────────────────────────────────────────────
\subsection{Apprendimento supervisionato}
% ─────────────────────────────────────────────────────────────

Il \textbf{learning from labels} (apprendimento supervisionato) è il
paradigma in cui ogni esempio del dataset è una coppia $(x, y)$ dove $x$ è
un vettore di feature e $y$ è l'etichetta vera (la ``risposta corretta'').
L'obiettivo del modello è imparare una funzione $f$ tale che $f(x) \approx y$
per esempi mai visti prima.

Dopo l'addestramento il modello deve \emph{generalizzare}: dare la risposta
giusta su input nuovi. Il modo più semplice per misurarlo è separare i
dati in due (o tre) sottoinsiemi disgiunti.

% ─────────────────────────────────────────────────────────────
\subsection{Train, Validation e Test set}
% ─────────────────────────────────────────────────────────────

\begin{itemize}
    \item \textbf{Training set} (tipicamente 60--80\% dei dati). Su questo
          il modello aggiusta i suoi parametri --- in una rete neurale, i
          pesi vengono modificati dalla backpropagation.
    \item \textbf{Validation set} (10--20\%). Non viene usato per
          aggiornare i pesi, ma per scegliere fra modelli/configurazioni
          (es. fra architetture diverse, learning rate diversi). Nel nostro
          progetto è il set su cui il GA misura la fitness di ogni
          individuo.
    \item \textbf{Test set} (10--20\%). Toccato solo alla fine, per
          stimare onestamente quanto il modello generalizza. Se lo si usa
          durante la ricerca, si introduce \textbf{data leakage}: il
          modello impara a fare bene su dati che dovrebbero essere segreti.
\end{itemize}

\begin{concetto}[separazione delle responsabilità tra i set]
Train, validation e test svolgono tre ruoli \emph{diversi e non
intercambiabili}:
\begin{itemize}
    \item Il train serve a \emph{imparare}.
    \item Il validation serve a \emph{decidere}.
    \item Il test serve a \emph{misurare} (una volta sola).
\end{itemize}
Confondere questi ruoli porta a stime gonfiate delle prestazioni.
\end{concetto}

\subsubsection{Holdout vs k-fold cross validation}

\textbf{Holdout} è il metodo più semplice: si divide il dataset una volta
sola in train/val/test. Veloce, ma il risultato dipende da come è stata
fatta la divisione. Funziona bene quando il dataset è abbastanza grande
da essere rappresentativo in ogni split.

\textbf{K-fold cross validation} divide il dataset in $k$ pezzi
(\emph{fold}) di uguale dimensione, addestra $k$ volte usando ogni volta
un fold diverso come validation, e media i risultati. Stima più affidabile,
ma costa $k$ volte di più. In questo progetto si usa l'holdout (sufficiente
data la mole dei dataset utilizzati e il costo già alto del GA).

\subsubsection{Stratificazione}

Quando si fa lo split è importante mantenere le proporzioni delle classi
in ciascun sottoinsieme. Senza stratificazione, su dataset piccoli può
capitare che una classe finisca quasi tutta nel training set e quasi
niente nel validation set --- il modello sembrerà funzionare benissimo
ma in realtà non ha mai visto certi casi.

% ─────────────────────────────────────────────────────────────
\subsection{Classificazione vs Regressione}
% ─────────────────────────────────────────────────────────────

Nel supervisionato si distinguono due tipi di problemi a seconda di cosa
è $y$:

\begin{itemize}
    \item \textbf{Classificazione} --- $y$ è una categoria discreta.
          \emph{``Questo fiore è setosa, versicolor o virginica?''}.
          La metrica naturale è l'\emph{accuracy} (\% di predizioni
          corrette). Sull'output si usa tipicamente \fnt{softmax} per
          ottenere una distribuzione di probabilità sulle classi.
    \item \textbf{Regressione} --- $y$ è un valore continuo.
          \emph{``Quanto costerà questa casa?''}. La metrica è MSE/RMSE
          (errore quadratico medio). Sull'output si usa una funzione
          lineare.
\end{itemize}

Il progetto è interamente di \emph{classificazione}: tutti e quattro i
dataset di test (Iris, Wine, Breast Cancer, Digits) chiedono di assegnare
ogni esempio a una categoria.

% ─────────────────────────────────────────────────────────────
\subsection{Iperparametri vs parametri}
% ─────────────────────────────────────────────────────────────

\begin{itemize}
    \item I \textbf{parametri} di un modello sono i valori che il modello
          \emph{impara dai dati} durante l'addestramento --- in una rete
          neurale, i pesi e i bias.
    \item Gli \textbf{iperparametri} sono le scelte di configurazione che
          fissiamo \emph{prima} dell'addestramento: learning rate, numero
          di epoche, dimensione della popolazione del GA, mutation rate,
          $\lambda$ della penalità di complessità, ecc. Anche
          l'\emph{architettura} stessa è un iperparametro --- ed è
          esattamente quella che il GA cerca di scegliere automaticamente.
\end{itemize}

\begin{concetto}[la struttura tipica di un esperimento di ML]
\begin{enumerate}
    \item Si fissano gli iperparametri (a mano o tramite ricerca automatica).
    \item Si addestra il modello sul training set facendo scendere il
          gradiente della loss.
    \item Si misurano le prestazioni sul validation set.
    \item Si ripetono i passi 1--3 con configurazioni diverse e si sceglie
          la migliore.
    \item Si valuta una sola volta la configurazione vincente sul test set
          per riportare il risultato finale.
\end{enumerate}
\end{concetto}

% ─────────────────────────────────────────────────────────────
\subsection{Overfitting e generalizzazione}
% ─────────────────────────────────────────────────────────────

Un modello in \textbf{overfitting} ha imparato a memoria i dati di
training invece di estrarne le regole generali. Si manifesta con un
\emph{accuracy alta sul training set ma bassa sul validation/test set}.
Le cause più comuni:
\begin{itemize}
    \item modello troppo grande rispetto alla quantità di dati (troppi
          parametri liberi);
    \item training troppo lungo (la rete si specializza sui dettagli);
    \item pattern accidentali nei dati di training (rumore).
\end{itemize}

L'\textbf{underfitting} è il problema opposto: il modello è troppo
semplice per cogliere la struttura del problema --- accuracy bassa sia
sul training che sul validation.

\begin{concetto}[bias/variance trade-off]
\textbf{Bias alto / variance bassa} = underfitting: il modello fa errori
sistematici perché è troppo rigido.
\textbf{Bias basso / variance alta} = overfitting: il modello cambia
molto al variare dei dati e impara il rumore.
La sweet spot è in mezzo, e spesso la si raggiunge regolarizzando il
modello --- nel nostro caso con la penalità $\lambda$ sulla complessità,
che incentiva architetture più snelle.
\end{concetto}

\begin{riassunto}
\begin{itemize}
    \item Il dataset si divide in train (impara), validation (decide), test (misura una sola volta).
    \item I parametri si imparano dai dati; gli iperparametri si fissano prima del training.
    \item L'architettura è un iperparametro --- ed è quello che il GA cerca automaticamente.
    \item L'overfitting si combatte con dataset adeguati, regolarizzazione o vincoli sulla complessità.
\end{itemize}
\end{riassunto}

\newpage

% =============================================================
\section{Reti Neurali Artificiali}
% =============================================================

% ─────────────────────────────────────────────────────────────
\subsection{Dal neurone biologico al neurone artificiale}
% ─────────────────────────────────────────────────────────────

Una rete neurale artificiale è un modello computazionale ispirato al
funzionamento del cervello biologico. L'unità elementare è il
\textbf{neurone}: riceve segnali in ingresso, li combina linearmente
tramite \emph{pesi} (le ``sinapsi''), aggiunge un \emph{bias} e applica
una funzione di attivazione non lineare per produrre l'uscita.

Per un singolo neurone $i$:
\begin{equation}
    a_i = f(\bm{w}_i^\top \bm{x} + b_i)
\end{equation}
dove $\bm{w}_i$ è il vettore dei pesi, $\bm{x}$ è il vettore di input,
$b_i$ il bias, $f$ la funzione di attivazione e $a_i$ l'uscita
(\emph{attivazione}).

Per un intero \textbf{layer} di neuroni si raggruppa in forma matriciale:
\begin{equation}
    \bm{a} = f(\bm{W}\bm{x} + \bm{b})
\end{equation}
con $\bm{W}$ matrice dei pesi (una riga per neurone) e $\bm{b}$ vettore dei
bias.

% ─────────────────────────────────────────────────────────────
\subsection{Topologia: input, hidden, output}
% ─────────────────────────────────────────────────────────────

Una rete \emph{feed-forward} (l'unica usata in questo progetto) è
strutturata in tre tipi di layer in cascata:

\begin{itemize}
    \item \textbf{Layer di input} --- riceve le feature del dataset, una
          per neurone. Non applica nessuna trasformazione, è un puro
          punto di ingresso.
    \item \textbf{Hidden layer(s)} --- uno o più layer intermedi. Ogni
          layer trasforma la rappresentazione dei dati applicando
          $\bm{W}\bm{x} + \bm{b}$ seguito da una funzione di attivazione
          non lineare. È qui che la rete impara feature di livello
          superiore.
    \item \textbf{Layer di output} --- produce la predizione finale. La
          funzione di attivazione dipende dal problema: \fnt{softmax}
          per classificazione multiclasse, \fnt{sigmoid} per
          classificazione binaria, lineare per regressione.
\end{itemize}

\textbf{Perché servono le funzioni di attivazione non lineari?} Senza di
esse una rete con $N$ layer è matematicamente equivalente a una rete con
$1$ layer: la composizione di trasformazioni lineari è lineare. La non
linearità è ciò che permette alla rete di rappresentare funzioni
arbitrariamente complesse (\emph{universal approximation theorem}).

% ─────────────────────────────────────────────────────────────
\subsection{Funzioni di attivazione}
% ─────────────────────────────────────────────────────────────

\begin{description}
    \item[ReLU] (Rectified Linear Unit) ---
        $f(x) = \max(0, x)$.
        Veloce, derivata semplice ($1$ se $x>0$, $0$ altrimenti). Funziona
        molto bene nei layer nascosti perché non satura per $x>0$.
        Problema: \emph{dying ReLU} --- se un neurone riceve sempre input
        negativi smette di apprendere (gradiente sempre $0$).
    \item[Leaky ReLU] ---
        $f(x) = x$ se $x>0$, altrimenti $\alpha x$ con $\alpha$ piccolo
        (es. $0.01$). Risolve il dying ReLU mantenendo un piccolo
        gradiente per $x<0$.
    \item[Sigmoid] ---
        $f(x) = \frac{1}{1+e^{-x}}$. Output in $(0,1)$.
        Usata storicamente nei layer nascosti, oggi sconsigliata perché
        \emph{satura} (per $|x|$ grande la derivata si avvicina a zero,
        causando il \emph{vanishing gradient}). Resta utile come
        attivazione di output per problemi binari.
    \item[Linear] ---
        $f(x) = x$. Identità. Usata in output per regressione, raramente
        nei layer nascosti.
    \item[Softmax] ---
        $f(x_i) = \frac{e^{x_i}}{\sum_j e^{x_j}}$. Trasforma un vettore in
        una distribuzione di probabilità (somma a $1$). È la scelta
        standard per l'output di una classificazione multiclasse: il
        neurone con probabilità più alta indica la classe predetta.
\end{description}

\begin{concetto}[funzioni di attivazione e GA]
Nel progetto le funzioni di attivazione disponibili sono ReLU, Leaky ReLU,
Sigmoid, Linear e Softmax. Per i layer nascosti il GA può scegliere
liberamente fra le prime quattro; sull'output la \fnt{softmax} è fissa
(tutti i task sono di classificazione multiclasse). La possibilità di
mescolare attivazioni diverse fra layer è una scelta progettuale: il GA
può combinare ReLU al primo layer e Leaky ReLU al secondo, ad esempio.
\end{concetto}

% ─────────────────────────────────────────────────────────────
\subsection{Forward step}
% ─────────────────────────────────────────────────────────────

Il \textbf{forward step} (o \emph{forward pass}) propaga un input
attraverso la rete per produrre la predizione. Per una rete con $L$
hidden layer:

\begin{equation}
    \bm{a}^{(0)} = \bm{x}, \qquad
    \bm{a}^{(\ell)} = f^{(\ell)}\!\left(\bm{W}^{(\ell)}\bm{a}^{(\ell-1)} + \bm{b}^{(\ell)}\right)
    \quad \ell = 1, \ldots, L+1
\end{equation}

L'output finale $\hat{\bm{y}} = \bm{a}^{(L+1)}$ è la predizione della rete.
Durante il training si confronta $\hat{\bm{y}}$ con l'etichetta vera $\bm{y}$
e si calcola la \emph{loss} (errore).

% ─────────────────────────────────────────────────────────────
\subsection{Loss functions}
% ─────────────────────────────────────────────────────────────

La \textbf{funzione di loss} misura quanto la predizione si discosta dalla
verità. La scelta dipende dal problema:

\begin{itemize}
    \item \textbf{Mean Squared Error (MSE)} --- per regressione:
        \[ \mathcal{L} = \frac{1}{n}\sum_i (\hat{y}_i - y_i)^2 \]
    \item \textbf{Binary Cross-Entropy} --- per classificazione binaria:
        \[ \mathcal{L} = -\bigl(y\log\hat{y} + (1-y)\log(1-\hat{y})\bigr) \]
    \item \textbf{Categorical Cross-Entropy} --- per classificazione
        multiclasse (l'unica usata in questo progetto):
        \[ \mathcal{L} = -\sum_i y_i\log\hat{y}_i \]
        dove $\bm{y}$ è in \emph{one-hot encoding} e $\hat{\bm{y}}$ è
        l'uscita della softmax.
\end{itemize}

\begin{concetto}[softmax + cross-entropy: la coppia che semplifica i conti]
Quando l'output è softmax e la loss è categorical cross-entropy, il
gradiente combinato si scrive in modo molto compatto:
\[
    \frac{\partial \mathcal{L}}{\partial \bm{z}^{(L+1)}} = \hat{\bm{y}} - \bm{y}
\]
dove $\bm{z}^{(L+1)}$ è il pre-attivazione dell'output. Questo risultato è
sfruttato nel codice per evitare di calcolare separatamente la derivata
della softmax (che è una matrice jacobiana) e quella della cross-entropy.
\end{concetto}

% ─────────────────────────────────────────────────────────────
\subsection{Backward step e Gradient Descent}
% ─────────────────────────────────────────────────────────────

Il \textbf{backward step} (o \emph{backpropagation}) è l'algoritmo che
calcola, per ogni peso $w$ della rete, la derivata parziale
$\frac{\partial \mathcal{L}}{\partial w}$ (il \emph{gradiente}). Una volta
ottenuto, si aggiornano i pesi nella direzione opposta al gradiente:

\begin{equation}
    w \;\leftarrow\; w - \eta \frac{\partial \mathcal{L}}{\partial w}
\end{equation}

dove $\eta$ è il \textbf{learning rate}. Questa è la \emph{discesa del
gradiente} (gradient descent).

\subsubsection{L'intuizione}

Il gradiente è la pendenza della funzione di loss rispetto al peso.
Spostarsi nella direzione opposta significa scendere lungo la pendenza,
verso il minimo della loss. Il learning rate $\eta$ controlla la
\emph{dimensione del passo}:

\begin{itemize}
    \item $\eta$ troppo \emph{piccolo} → la convergenza è lentissima, in
          $N$ epoche fissate la rete non arriva al minimo.
    \item $\eta$ troppo \emph{grande} → i passi sono così ampi che la
          rete ``salta'' oltre il minimo, oscilla o diverge.
    \item $\eta$ \emph{giusto} → discesa rapida e stabile fino al minimo.
\end{itemize}

\subsubsection{Backpropagation: l'idea}

Il problema di calcolare $\frac{\partial \mathcal{L}}{\partial w}$ per
\emph{tutti} i pesi della rete è che la loss dipende dal peso $w$ in modo
indiretto: $w$ contribuisce all'attivazione di un neurone, che contribuisce
a quella del layer successivo, e così via fino all'output. Bisogna quindi
applicare la \textbf{regola della catena} ripetutamente.

L'osservazione chiave: i calcoli del layer $\ell$ riusano risultati del
layer $\ell+1$. Se si parte dall'output e si va all'indietro layer per
layer, si calcola tutto in un singolo \emph{forward pass} seguito da un
singolo \emph{backward pass} --- ecco perché si chiama
\emph{back-propagation}.

\begin{algorithm}[H]
\caption{Backpropagation (schema concettuale)}
\begin{algorithmic}[1]
\Require input $\bm{x}$, etichetta $\bm{y}$, rete con $L$ hidden layer
\Ensure aggiornamento di tutti i pesi $\bm{W}^{(\ell)}$ e bias $\bm{b}^{(\ell)}$
\State \textbf{Forward:} calcola $\bm{a}^{(\ell)}$ per $\ell = 1, \dots, L+1$ e memorizza tutto
\State Calcola la loss $\mathcal{L}(\hat{\bm{y}}, \bm{y})$
\State \textbf{Backward:} calcola $\bm{\delta}^{(L+1)} = \nabla_{\bm{z}^{(L+1)}} \mathcal{L}$
\For{$\ell = L$ downto $1$}
    \State $\bm{\delta}^{(\ell)} \gets (\bm{W}^{(\ell+1)})^\top \bm{\delta}^{(\ell+1)} \odot f'^{(\ell)}\!(\bm{z}^{(\ell)})$
\EndFor
\For{$\ell = 1$ to $L+1$}
    \State $\bm{W}^{(\ell)} \gets \bm{W}^{(\ell)} - \eta\, \bm{\delta}^{(\ell)}(\bm{a}^{(\ell-1)})^\top$
    \State $\bm{b}^{(\ell)} \gets \bm{b}^{(\ell)} - \eta\, \bm{\delta}^{(\ell)}$
\EndFor
\end{algorithmic}
\end{algorithm}

Le quantità $\bm{\delta}^{(\ell)}$ sono il ``gradiente locale'' di ogni
layer, calcolato in modo ricorsivo dall'output verso l'input.

\subsubsection{SGD vs batch vs mini-batch}

Esistono tre modalità di gradient descent a seconda di quanti esempi si
usano per stimare il gradiente prima di aggiornare i pesi:
\begin{itemize}
    \item \textbf{Batch GD} --- usa tutti gli esempi del training set
          per ogni aggiornamento. Stima accurata, ma lenta e costosa.
    \item \textbf{Stochastic GD (SGD)} --- usa un solo esempio alla volta.
          Stima rumorosa ma il rumore stesso aiuta a uscire da minimi
          locali. È la modalità usata in questo progetto.
    \item \textbf{Mini-batch GD} --- compromesso: $32, 64, 128$ esempi
          alla volta. Standard nei framework moderni.
\end{itemize}

% ─────────────────────────────────────────────────────────────
\subsection{Inizializzazione dei pesi}
% ─────────────────────────────────────────────────────────────

I pesi non possono essere inizializzati a zero (tutti i neuroni
diventerebbero identici e imparerebbero la stessa cosa --- problema della
\emph{symmetry breaking}) né a valori troppo grandi (gradient explosion).
Le strategie più usate sono:

\begin{itemize}
    \item \textbf{Inizializzazione di He} --- pesi $\sim \mathcal{N}(0,
          \sigma^2)$ con $\sigma^2 = 2/n_{\text{in}}$. Ottimizzata per
          ReLU/Leaky ReLU.
    \item \textbf{Inizializzazione di Xavier (o Glorot)} --- $\sigma^2 =
          1/n_{\text{in}}$. Ottimizzata per sigmoid e tanh.
\end{itemize}

In questo progetto si applica He per ReLU e Leaky ReLU, Xavier per le
altre funzioni. È un dettaglio importante: senza inizializzazione
ragionata, reti più profonde tendono al \emph{vanishing} o
\emph{exploding gradient}.

\begin{concetto}[gradient clipping]
Anche con buona inizializzazione, durante il training i gradienti possono
``esplodere'' a valori giganteschi a causa di catene di moltiplicazioni.
Il \emph{gradient clipping} è una protezione: se la norma del gradiente
supera una soglia, si scala il gradiente per riportarlo entro la soglia.
È implementato nella backpropagation del progetto.
\end{concetto}

\begin{riassunto}
\begin{itemize}
    \item Una rete neurale è una composizione di trasformazioni lineari $\bm{W}\bm{x}+\bm{b}$ alternate a non linearità $f$.
    \item Le non linearità sono \emph{necessarie}: senza, la rete è equivalente a un singolo layer lineare.
    \item Il forward step calcola la predizione; la loss misura l'errore; il backward step calcola il gradiente per scenderlo.
    \item Il learning rate $\eta$ è l'iperparametro più critico della backpropagation: controlla la dimensione del passo.
    \item Backpropagation = regola della catena applicata in modo efficiente partendo dall'output.
\end{itemize}
\end{riassunto}

\newpage

% =============================================================
\section{Algoritmi Genetici}
% =============================================================

% ─────────────────────────────────────────────────────────────
\subsection{La metafora biologica}
% ─────────────────────────────────────────────────────────────

Gli \textbf{algoritmi genetici} (GA) sono una classe di algoritmi
evolutivi ispirati alla selezione naturale di Darwin. L'idea fondante è
che soluzioni a un problema possono ``evolvere'' nel tempo se sottoposte
ai meccanismi della selezione naturale: i ``più adatti'' sopravvivono e
si riproducono, i ``meno adatti'' tendono a scomparire.

Nel quadro di Russell \& Norvig, un GA può essere visto come una variante
della \emph{ricerca locale stocastica}: si mantiene non un singolo stato
candidato ma un'\emph{intera popolazione} di stati, e si esplora lo spazio
combinando coppie di stati invece di muoversi a piccoli passi da uno
stato singolo.

\begin{concetto}[un GA è caratterizzato da queste scelte]
\begin{itemize}
    \item Dimensione della \emph{popolazione}.
    \item \emph{Rappresentazione} di ogni individuo (stringa, lista, vettore\ldots).
    \item Numero $\rho$ di genitori per ogni figlio (tipicamente $\rho=2$).
    \item Procedura di \emph{selezione} dei genitori (a torneo, a roulette, \ldots).
    \item Procedura di \emph{ricombinazione} (crossover).
    \item \emph{Tasso di mutazione}.
    \item Strategia di formazione della nuova generazione (con o senza elitismo).
\end{itemize}
\end{concetto}

% ─────────────────────────────────────────────────────────────
\subsection{Componenti fondamentali}
% ─────────────────────────────────────────────────────────────

\begin{description}
    \item[Individuo (cromosoma)] Una soluzione candidata al problema.
        Nel progetto è una \texttt{lista di tuple}, dove ogni tupla
        \texttt{(n\_neuroni, funzione)} rappresenta un hidden layer.
    \item[Gene] Elemento atomico del cromosoma. Nel progetto, una singola
        tupla \texttt{(n\_neuroni, funzione)}.
    \item[Popolazione] Insieme di tutti gli individui in una generazione.
    \item[Generazione] Una iterazione completa del ciclo evolutivo.
    \item[Fitness] Funzione che misura la qualità di un individuo. È il
        ``segnale evolutivo'' che guida la ricerca: più alta è meglio.
\end{description}

% ─────────────────────────────────────────────────────────────
\subsection{Il ciclo evolutivo}
% ─────────────────────────────────────────────────────────────

Lo schema generale di un GA è il seguente:

\begin{algorithm}[H]
\caption{Algoritmo Genetico (schema generale)}
\begin{algorithmic}[1]
\State Genera una popolazione iniziale casuale di $N$ individui
\Repeat
    \State Valuta la fitness di ogni individuo
    \State \textbf{Elitismo:} copia il migliore nella nuova popolazione
    \While{nuova popolazione < $N$}
        \State \emph{Selezione:} scegli due genitori dalla popolazione corrente
        \State \emph{Crossover:} combina i due genitori per creare un figlio
        \State \emph{Mutazione:} con probabilità $p_m$ modifica il figlio
        \State Aggiungi il figlio alla nuova popolazione
    \EndWhile
    \State popolazione $\gets$ nuova popolazione
\Until{condizione di terminazione (numero di generazioni o fitness target)}
\State \textbf{return} l'individuo con fitness più alta mai trovato
\end{algorithmic}
\end{algorithm}

% ─────────────────────────────────────────────────────────────
\subsection{Selezione: torneo}
% ─────────────────────────────────────────────────────────────

Esistono diverse procedure di selezione, ognuna con un diverso modo di
bilanciare due esigenze opposte:

\begin{itemize}
    \item \textbf{Sfruttamento (\emph{exploitation})} --- i migliori
          individui devono avere più probabilità di riprodursi.
    \item \textbf{Esplorazione (\emph{exploration})} --- anche individui
          mediocri devono avere qualche chance, per mantenere diversità
          genetica e ridurre il rischio di \emph{convergenza prematura}
          su un ottimo locale.
\end{itemize}

Le tre procedure più comuni:

\begin{itemize}
    \item \textbf{Roulette wheel} --- ogni individuo ha probabilità di
          essere scelto proporzionale alla sua fitness. Semplice ma sensibile
          alla scala della fitness (un individuo molto migliore degli altri
          domina subito).
    \item \textbf{Selezione a torneo} --- (usata nel progetto). Si estraggono
          $k$ individui a caso dalla popolazione e si sceglie il migliore tra
          loro come genitore. Non sensibile alla scala assoluta della fitness,
          solo all'\emph{ordinamento}: un individuo con fitness $0.5$ ha la
          stessa probabilità di vincere su un avversario con fitness $0.4$
          che su uno con fitness $0.1$.
    \item \textbf{Selezione per rango} --- gli individui sono ordinati per
          fitness, le probabilità sono assegnate in base al rango (non al
          valore). Anche questa è insensibile alla scala.
\end{itemize}

\begin{domanda}
\textit{``Perché selezione a torneo invece di scegliere sempre i migliori?''}\\
Se selezionassi sempre i migliori, dopo poche generazioni l'intera
popolazione diventerebbe quasi identica e il GA smetterebbe di esplorare,
restando incastrato in un ottimo locale. Con il torneo $k=3$, un individuo
forte ha alta probabilità di vincere ma uno mediocre ha ancora chance se
gli capita un torneo con avversari più deboli. La taglia $k$ controlla il
bilanciamento: $k$ grande $\Rightarrow$ più pressione selettiva, $k$
piccolo $\Rightarrow$ più esplorazione.
\end{domanda}

% ─────────────────────────────────────────────────────────────
\subsection{Crossover}
% ─────────────────────────────────────────────────────────────

Il \textbf{crossover} (o \emph{ricombinazione}) combina due genitori per
generare un figlio. L'idea: se entrambi i genitori sono buoni, è probabile
che combinazioni di parti di ciascuno producano figli ancora migliori.

La forma classica per stringhe di lunghezza fissa è il \emph{single-point
crossover}: si sceglie un punto di taglio casuale $c$, il figlio prende
le posizioni $1\dots c$ dal genitore 1 e $c+1\dots n$ dal genitore 2.

\begin{verbatim}
genitore_1: [g1 g1 g1 | g1 g1 g1]   <- punto di taglio c
genitore_2: [g2 g2 g2 | g2 g2 g2]
figlio_1:   [g1 g1 g1 | g2 g2 g2]
\end{verbatim}

Nel progetto i cromosomi hanno \emph{lunghezza variabile} (architetture
con numero di layer diverso), quindi serve una variante: si scelgono due
punti di taglio indipendenti, $t_1$ per il primo genitore e $t_2$ per il
secondo, e si concatena.

% ─────────────────────────────────────────────────────────────
\subsection{Mutazione}
% ─────────────────────────────────────────────────────────────

La \textbf{mutazione} introduce variazioni casuali nel cromosoma. È
fondamentale per due motivi:
\begin{itemize}
    \item Mantiene \emph{diversità genetica} nella popolazione, evitando
          la convergenza prematura.
    \item Permette di esplorare zone dello spazio delle soluzioni che il
          crossover da solo non potrebbe raggiungere (il crossover ricombina
          ma non crea ``materiale genetico'' davvero nuovo).
\end{itemize}

Il \emph{mutation rate} $p_m$ è la probabilità che ciascun gene venga
modificato. Valori tipici: $0.05$--$0.30$.

\subsubsection{I tre tipi di mutazione del progetto}

Ogni gene è una tupla \texttt{(n\_neuroni, funzione)}. Quando un gene viene
mutato, una delle tre seguenti operazioni viene scelta con uguale
probabilità:
\begin{enumerate}
    \item \textbf{Mutazione del numero di neuroni} --- si sostituisce
          \texttt{n\_neuroni} con un nuovo valore casuale entro il range
          ammesso, mantenendo la stessa funzione.
    \item \textbf{Mutazione della funzione} --- si sostituisce la funzione
          di attivazione con un'altra del pool, mantenendo lo stesso numero
          di neuroni.
    \item \textbf{Mutazione strutturale} --- si aggiunge un nuovo layer in
          fondo o si rimuove uno dei layer esistenti (a meno che il
          cromosoma non abbia un solo layer, nel qual caso non si rimuove).
\end{enumerate}

I primi due tipi raffinano un'architettura esistente; il terzo permette di
esplorare \emph{topologie} diverse (numero di layer variabile).

% ─────────────────────────────────────────────────────────────
\subsection{Elitismo}
% ─────────────────────────────────────────────────────────────

L'\textbf{elitismo} è una tecnica semplice e potente: il miglior individuo
della generazione corrente viene copiato \emph{integralmente} nella
generazione successiva, senza subire crossover o mutazione. Questo
garantisce che la \emph{best fitness} della popolazione non possa mai
peggiorare da una generazione all'altra.

\begin{concetto}[perché elitismo]
Senza elitismo, il miglior individuo trovato finora potrebbe essere
``perso'' a causa della casualità di crossover/mutazione. Con elitismo, la
ricerca è \emph{monotona}: una volta trovata un'architettura buona, non
viene più persa. Lo svantaggio è una leggera riduzione della diversità,
ma in pratica copiare 1 individuo su 20 ha effetto trascurabile.
\end{concetto}

% ─────────────────────────────────────────────────────────────
\subsection{GA come variante della ricerca locale}
% ─────────────────────────────────────────────────────────────

Vale la pena collocare il GA nel quadro più ampio degli algoritmi di
ricerca locale visto a lezione (slide ``Ricerca in ambienti complessi'').
Gli algoritmi di ricerca locale lavorano nel \emph{problem space} senza
mantenere il percorso, e cercano una soluzione migliorando un singolo
stato (o un piccolo insieme di stati) tramite mosse locali.

\begin{itemize}
    \item \textbf{Hill climbing} --- a ogni passo va al successore con
          valore migliore. Veloce ma rischia ottimi locali.
    \item \textbf{Simulated annealing} --- accetta anche mosse peggiorative
          con probabilità decrescente nel tempo. Permette di uscire da
          ottimi locali.
    \item \textbf{Local beam search} --- mantiene $k$ stati in parallelo,
          a ogni passo li espande tutti e tiene i migliori $k$ successori.
    \item \textbf{Algoritmi genetici} --- come il beam search ma con
          ricombinazione fra stati: i figli sono generati \emph{combinando}
          due stati genitore, non da successori di un singolo stato.
\end{itemize}

\begin{domanda}
\textit{``Cosa distingue un GA da un beam search?''}\\
Il beam search mantiene $k$ stati e a ogni passo genera i successori di
ciascuno indipendentemente. Il GA aggiunge il \emph{crossover}: i nuovi
candidati possono combinare informazioni da più stati genitori. Questo
rende il GA capace di mescolare ``ingredienti buoni'' provenienti da
soluzioni diverse, cosa che il beam search non fa.
\end{domanda}

\begin{riassunto}
\begin{itemize}
    \item Un GA è un algoritmo di ricerca evolutiva che lavora su una popolazione di soluzioni candidate.
    \item Le componenti fondamentali sono: rappresentazione, fitness, selezione, crossover, mutazione.
    \item Le scelte di progettazione del GA bilanciano \emph{esplorazione} (mutazione, popolazione grande) e \emph{sfruttamento} (selezione, elitismo).
    \item Il GA si colloca nella famiglia delle ricerche locali stocastiche, distinto da hill climbing e simulated annealing per la presenza della popolazione e del crossover.
\end{itemize}
\end{riassunto}

\newpage

% =============================================================
\section{Implementazione del progetto}
% =============================================================

% ─────────────────────────────────────────────────────────────
\subsection{Architettura del sistema}
% ─────────────────────────────────────────────────────────────

Il sistema è composto da due moduli principali che collaborano:

\begin{itemize}
    \item Il modulo \texttt{neural\_network} contiene la classe
          \texttt{neural\_network} (motore feedforward + backpropagation)
          e le funzioni di attivazione.
    \item Il modulo \texttt{genetic\_algorithm} contiene la classe
          \texttt{GeneticAlgorithm} che orchestra il ciclo evolutivo e
          chiama la rete come ``oracolo'' per valutare le architetture
          candidate.
\end{itemize}

Il flusso di alto livello è:

\begin{verbatim}
     +---------------------------------------------+
     | Carica dataset -> train / val / test split  |
     +---------------------------+-----------------+
                                 |
                                 v
     +---------------------------------------------+
     | Inizializza popolazione casuale di N reti   |
     +---------------------------+-----------------+
                                 |
                                 v
     +---------------------------------------------+
     | Per ogni generazione:                       |
     |   1. valuta fitness di ogni individuo       |
     |      (per ogni individuo: training BP +     |
     |       accuracy su val, ripetuto K volte)    |
     |   2. salva il migliore (elitismo)           |
     |   3. genera nuova popolazione tramite       |
     |      selezione + crossover + mutazione      |
     +---------------------------+-----------------+
                                 |
                                 v
     +---------------------------------------------+
     | Riaddestra il migliore e valuta sul TEST    |
     | set (stima onesta delle prestazioni)        |
     +---------------------------------------------+
\end{verbatim}

% ─────────────────────────────────────────────────────────────
\subsection{Rappresentazione del cromosoma}
% ─────────────────────────────────────────────────────────────

Un individuo è una lista Python di tuple, una per ogni hidden layer:

\begin{verbatim}
[(12, relu), (6, sigmoid)]              # 2 hidden layer
[(32, relu), (16, relu), (8, relu)]     # 3 hidden layer
[(8, linear)]                           # 1 hidden layer
\end{verbatim}

\begin{itemize}
    \item Il \emph{numero} di tuple = numero di hidden layer (può variare
          da individuo a individuo).
    \item Per ogni tupla: \texttt{n\_neuroni} (intero in un range fissato,
          es. $4$--$64$) e \texttt{funzione} (scelta fra ReLU, Leaky ReLU,
          Sigmoid, Linear).
    \item Il layer di input e quello di output sono \emph{determinati dal
          dataset} (numero di feature in input, numero di classi in
          output) e non fanno parte del cromosoma.
    \item La funzione di output è fissa a \fnt{softmax}, perché tutti i
          dataset usati sono di classificazione multiclasse.
\end{itemize}

\begin{domanda}
\textit{``I pesi della rete fanno parte del cromosoma?''}\\
No. Il cromosoma rappresenta solo l'\emph{architettura} (struttura
discreta). I pesi sono inizializzati casualmente \emph{ogni volta} che si
costruisce una rete per valutarla. La metafora biologica è precisa: i
geni trasmettono le istruzioni per costruire l'organismo, non l'organismo
già formato. L'architettura è il gene, i pesi sono le caratteristiche
fisiche sviluppate durante il training.
\end{domanda}

% ─────────────────────────────────────────────────────────────
\subsection{Funzione di fitness}
% ─────────────────────────────────────────────────────────────

La fitness misura quanto è ``buona'' un'architettura. Questa è la
formulazione finale usata nel progetto:

\begin{equation}
    \text{fitness} = \text{accuracy} - \lambda \cdot \log_{10}(N_{\text{params}})
\end{equation}

dove $\text{accuracy}$ è la percentuale di predizioni corrette sul
validation set (mediata su $K$ run indipendenti), $N_{\text{params}}$ è
il numero totale di pesi e bias dell'architettura, e $\lambda \geq 0$ è
il coefficiente di penalità di complessità.

\subsubsection{Perché penalizzare la complessità}

Senza penalità ($\lambda = 0$) il GA tende a far crescere le architetture
generazione dopo generazione, perché reti più grandi tendono ad avere
accuracy leggermente superiore (maggiore capacità di adattarsi ai dati).
Questo causa due problemi:
\begin{itemize}
    \item Reti grandi sono più soggette a \emph{overfitting} sul validation
          set (gap fra accuracy val e accuracy test maggiore).
    \item Reti grandi sono inutilmente costose da addestrare e usare.
\end{itemize}
La penalità $\lambda$ è una forma di \emph{regolarizzazione}: scoraggia la
complessità per scelta progettuale.

\subsubsection{Perché logaritmica e non lineare}

Una prima versione del progetto usava penalità lineare, $\lambda \cdot
N_{\text{params}}$. Funziona, ma il valore di $\lambda$ utile dipende
fortemente dalla scala delle reti generate: una stessa $\lambda$ non
penalizza praticamente nulla una rete piccola e annienta una rete grande.

Con penalità logaritmica, il valore della penalità cresce in modo molto
più lento: passare da $100$ a $10\,000$ parametri (un fattore $\times 100$)
fa salire il termine di soli $2$ punti ($\log_{10}(100)=2$,
$\log_{10}(10000)=4$). Questo rende $\lambda$ stabile e indipendente
dalla scala assoluta dei parametri.

\subsubsection{Perché K valutazioni}

I pesi della rete sono inizializzati casualmente a ogni valutazione, e
con SGD a campioni casuali il training è di per sé stocastico.
Conseguenza: una stessa architettura, addestrata due volte di seguito,
può ottenere accuracy diverse anche di alcuni punti percentuali. Questo
\emph{rumore di valutazione} può rendere il segnale di fitness
inaffidabile e portare il GA a selezionare individui ``fortunati''
piuttosto che individui realmente buoni.

La soluzione è ripetere la valutazione $K$ volte e mediare:
\[
    \text{accuracy} = \frac{1}{K}\sum_{k=1}^{K} \text{accuracy}_k
\]

Con $K=3$ la deviazione standard si dimezza rispetto a $K=1$ (vedi
sezione esperimenti). Aumentare $K$ oltre $3$ riduce ulteriormente la
varianza ma triplica o quintuplica il costo computazionale, con beneficio
marginale.

\begin{algorithm}[H]
\caption{Valutazione della fitness di un individuo}
\begin{algorithmic}[1]
\Require cromosoma $c$ (architettura), training set, validation set
\Ensure (fitness, accuracy)
\State $\text{sum\_acc} \gets 0$
\For{$k = 1$ to $K$}
    \State Costruisci la rete $\mathcal{N}$ con architettura $c$ e pesi inizializzati a caso
    \For{$\text{epoch} = 1$ to $E$}
        \State Estrai un campione casuale $(x, y)$ dal training set
        \State Esegui un passo di backpropagation su $(x, y)$
    \EndFor
    \State $\text{correct} \gets$ numero di predizioni corrette di $\mathcal{N}$ sul validation set
    \State $\text{acc}_k \gets \text{correct} \;/\; |\text{validation set}|$
    \State $\text{sum\_acc} \gets \text{sum\_acc} + \text{acc}_k$
\EndFor
\State $\text{accuracy} \gets \text{sum\_acc} / K$
\State $\text{fitness} \gets \text{accuracy} - \lambda \cdot \log_{10}(\#\text{params}(c))$
\State \textbf{return} $(\text{fitness}, \text{accuracy})$
\end{algorithmic}
\end{algorithm}

\textbf{Nota:} fitness e accuracy vengono restituite \emph{separatamente}.
Il GA usa la fitness per la selezione, ma nei grafici si mostra
l'accuracy pura per leggibilità.

% ─────────────────────────────────────────────────────────────
\subsection{Operatori genetici (pseudocodice)}
% ─────────────────────────────────────────────────────────────

\subsubsection{Selezione a torneo}

\begin{algorithm}[H]
\caption{\fnt{selezione\_torneo}(popolazione, fitness, $k$)}
\begin{algorithmic}[1]
\State Estrai casualmente $k$ indici dalla popolazione
\State Ordina i $k$ individui per fitness decrescente
\State \textbf{return} l'individuo con fitness più alta tra i $k$
\end{algorithmic}
\end{algorithm}

Viene chiamata \emph{due volte} per ogni figlio da generare (una per
ogni genitore). Il parametro $k$ (qui chiamato \texttt{tournament\_size})
controlla la pressione selettiva.

\subsubsection{Crossover a doppio punto di taglio}

\begin{algorithm}[H]
\caption{\fnt{crossover}(genitore$_1$, genitore$_2$)}
\begin{algorithmic}[1]
\State $t_1 \gets$ punto di taglio casuale in $[0, |\text{gen}_1|]$
\State $t_2 \gets$ punto di taglio casuale in $[0, |\text{gen}_2|]$
\State $\text{figlio} \gets \text{gen}_1[\,:t_1\,] + \text{gen}_2[\,t_2:\,]$
\If{$\text{figlio}$ è vuoto} \Comment{caso degenere}
    \State $\text{figlio} \gets [\,$un gene casuale dal pool dei due genitori$\,]$
\EndIf
\State \textbf{return} $\text{figlio}$
\end{algorithmic}
\end{algorithm}

\textbf{Perché due punti di taglio diversi e non uno solo?} Perché i due
genitori possono avere lunghezze diverse: con un singolo $t$ si
rischierebbero combinazioni nonsensiche o sempre della stessa lunghezza.
Con $t_1, t_2$ indipendenti il figlio può avere lunghezza qualsiasi tra
$0$ e $|gen_1| + |gen_2|$, esplorando bene lo spazio delle topologie.

\subsubsection{Mutazione}

\begin{algorithm}[H]
\caption{\fnt{mutazione}(individuo, $p_m$)}
\begin{algorithmic}[1]
\State $\text{nuovo} \gets$ copia di individuo
\For{ogni gene $(n, f)$ in $\text{nuovo}$}
    \If{rand $< p_m$}
        \State scegli casualmente uno dei tre tipi:
        \begin{itemize}
            \item[\textsc{Tipo 1}] sostituisci $n$ con un nuovo valore casuale (mantieni $f$)
            \item[\textsc{Tipo 2}] sostituisci $f$ con un'altra funzione casuale (mantieni $n$)
            \item[\textsc{Tipo 3}] aggiungi un nuovo gene casuale o rimuovi questo gene (se $|\text{nuovo}|>1$)
        \end{itemize}
    \EndIf
\EndFor
\State \textbf{return} $\text{nuovo}$
\end{algorithmic}
\end{algorithm}

% ─────────────────────────────────────────────────────────────
\subsection{Il ciclo principale (run)}
% ─────────────────────────────────────────────────────────────

\begin{algorithm}[H]
\caption{\fnt{GeneticAlgorithm.run()}}
\begin{algorithmic}[1]
\State $P \gets$ popolazione iniziale di $N$ individui casuali
\State Inizializza storia\_best\_fit, storia\_best\_acc, storia\_mean\_acc
\State $\text{best\_global} \gets$ \texttt{None}
\For{$g = 1$ to generazioni}
    \State Per ogni $\text{ind} \in P$, calcola $(\text{fit}, \text{acc}) = \fnt{fitness}(\text{ind})$
    \State Aggiorna best\_global se trovato un migliore
    \State Salva statistiche della generazione $g$ nelle storie
    \State $P_{\text{new}} \gets [\,$ best della generazione $\,]$ \Comment{elitismo}
    \While{$|P_{\text{new}}| < N$}
        \State $\text{p}_1 \gets \fnt{selezione\_torneo}(P, k)$
        \State $\text{p}_2 \gets \fnt{selezione\_torneo}(P, k)$
        \State $\text{f} \gets \fnt{crossover}(\text{p}_1, \text{p}_2)$
        \State $\text{f} \gets \fnt{mutazione}(\text{f}, p_m)$
        \State Aggiungi $\text{f}$ a $P_{\text{new}}$
    \EndWhile
    \State $P \gets P_{\text{new}}$
\EndFor
\State \textbf{return} (best\_global, storie)
\end{algorithmic}
\end{algorithm}

\begin{concetto}[parallelizzazione delle valutazioni]
La valutazione della fitness di ogni individuo è la fase più costosa
dell'algoritmo (richiede $K$ training completi). Gli individui di una
stessa generazione sono indipendenti fra loro, quindi possono essere
valutati \emph{in parallelo}. Nel progetto si usa
\texttt{ProcessPoolExecutor} di Python per distribuire le valutazioni su
più CPU.
\end{concetto}

% ─────────────────────────────────────────────────────────────
\subsection{Preprocessing del dataset}
% ─────────────────────────────────────────────────────────────

Tre operazioni di preprocessing standard, in questo ordine:

\subsubsection{Split train / validation / test}

Il dataset viene diviso in tre sottoinsiemi disgiunti, tipicamente con
proporzioni $60\%/20\%/20\%$. Lo split è \textbf{stratificato}: la
proporzione delle classi è preservata in ogni sottoinsieme. Il GA
valuta la fitness sul \emph{validation}; il \emph{test} viene toccato
solo alla fine, una volta sola, per riportare il risultato finale.

\subsubsection{Normalizzazione (StandardScaler)}

Le feature vengono normalizzate per avere media $0$ e deviazione
standard $1$:
\[ x' = \frac{x - \mu}{\sigma} \]
Questo aiuta la backpropagation perché tutte le feature operano sulla
stessa scala (altrimenti pesi diversi devono compensare scale diverse).

\begin{attenzione}
$\mu$ e $\sigma$ vanno calcolati \emph{solo sul training set}. Si applicano
poi gli stessi valori (senza ricalcolarli) al validation e al test set.
Calcolarli sull'intero dataset è una forma di \emph{data leakage}:
informazione del test set si ``infiltra'' nel training, gonfiando le
stime di performance. Questo è oggetto del Test 9B.
\end{attenzione}

\subsubsection{One-hot encoding}

Le etichette di classe (interi $0, 1, 2, \ldots$) vengono convertite in
vettori binari con un solo $1$ nella posizione corrispondente alla classe:
\[
    y = 2 \;\rightarrow\; \bm{y} = [0,\, 0,\, 1,\, 0,\, \ldots,\, 0]
\]
Questo è necessario perché la rete ha tanti neuroni di output quante sono
le classi, e la softmax produce una distribuzione di probabilità su tutte
le classi.

% ─────────────────────────────────────────────────────────────
\subsection{La baseline di confronto}
% ─────────────────────────────────────────────────────────────

Per dare senso ai numeri prodotti dal GA, ogni esperimento confronta il
risultato con una \textbf{baseline}: una rete con architettura
\emph{fissa} addestrata con backpropagation classica, senza l'aiuto del
GA. La baseline utilizzata è $[(8,\text{leaky\_relu}), (8,\text{leaky\_relu})]$.

\begin{itemize}
    \item Se il GA non batte la baseline, l'intero costo computazionale
          della ricerca evolutiva non è giustificato.
    \item Se il GA la batte, il valore aggiunto della ricerca automatica
          dell'architettura è dimostrato concretamente.
\end{itemize}

Per equità, la baseline è addestrata per lo stesso numero di epoche
totali usato per valutare un singolo individuo del GA: $K \cdot E$ epoche
(es. $3 \cdot 500 = 1500$).

Il numero di parametri della baseline varia a seconda del dataset (per via
del diverso numero di feature in input e di classi in output):

\begin{center}
\begin{tabular}{lccr}
    \toprule
    \rowcolor{headerblue}
    \textcolor{white}{\textbf{dataset}} &
    \textcolor{white}{\textbf{feat.~in}} &
    \textcolor{white}{\textbf{classi out}} &
    \textcolor{white}{\textbf{parametri baseline}} \\
    \midrule
    iris           &  4 &  3 & 139 \\
    \rowcolor{rowgray}
    wine           & 13 &  3 & 211 \\
    breast\_cancer & 30 &  2 & 338 \\
    \rowcolor{rowgray}
    digits         & 64 & 10 & 682 \\
    \bottomrule
\end{tabular}
\end{center}

Nei test di sensibilità (Test 1--7) il dataset è sempre \texttt{digits},
quindi il confronto è sempre rispetto a \textbf{682 parametri}.

\begin{riassunto}
\begin{itemize}
    \item Il sistema è composto dal modulo \texttt{neural\_network} (motore) e \texttt{genetic\_algorithm} (orchestratore).
    \item Il cromosoma è una lista di tuple \texttt{(n\_neuroni, funzione)}, una per hidden layer.
    \item La fitness è $\text{accuracy} - \lambda \cdot \log_{10}(N_{\text{params}})$, mediata su $K$ run.
    \item Crossover a doppio punto di taglio (per gestire lunghezze variabili), mutazione di 3 tipi (numero, funzione, struttura), elitismo.
    \item Lo split è stratificato e la normalizzazione \emph{si calcola solo sul training set} per evitare data leakage.
    \item La baseline è una rete fissa che permette di misurare se il GA aggiunge valore.
\end{itemize}
\end{riassunto}

\newpage

% =============================================================
\section{Concetti chiave trasversali}
% =============================================================

Questa sezione raccoglie alcuni concetti che attraversano il progetto e
che molto probabilmente verranno chiesti all'orale. Ognuno è collegato
all'una o all'altra delle scelte progettuali.

% ─────────────────────────────────────────────────────────────
\subsection{Overfitting nel contesto del progetto}
% ─────────────────────────────────────────────────────────────

\label{ssec:overfitting}

Il GA, lasciato senza vincoli, ha una naturale tendenza a far crescere
le architetture. Più una rete è grande, più parametri liberi ha, e più
è capace di adattarsi (anche \emph{troppo}) ai dati. Sul training set
questo si traduce in accuracy più alta; sul validation set anche, ma in
misura minore; sul test set spesso \emph{peggio}.

\begin{attenzione}
\textbf{Troppi parametri $\Rightarrow$ maggiore rischio di overfitting.}\\
Utilizzare una rete con molti più parametri di quanti ne servano
realmente \emph{aumenta} la probabilità di overfitting. Ecco perché:
\begin{itemize}
    \item Ogni parametro è un grado di libertà in più che il modello può
          usare per adattarsi ai dati.
    \item Se i parametri sono più del necessario, quelli ``in eccesso''
          non hanno pattern reali da imparare: finiscono per memorizzare
          il rumore specifico del training set.
    \item Questo si traduce in alta accuracy sul training ma
          generalizzazione peggiore sul test --- la definizione classica
          di overfitting.
\end{itemize}
\textbf{Esempio concreto (Test 8 --- Wine):} la baseline usa $211$
parametri, il GA solo $122$, entrambi raggiungono $100\%$. I $89$
parametri in eccesso della baseline non portano nessun beneficio di
accuracy, ma su un dataset più rumoroso o con meno campioni sarebbero
liberi di memorizzare pattern spurii. La rete GA, più compatta, avrebbe
un rischio di overfitting strutturalmente inferiore.
\end{attenzione}

\begin{concetto}[overfitting sul validation set]
È una forma sottile di overfitting: il GA non sta facendo overfitting
sul \emph{training} (la backpropagation lo gestisce), sta facendo
overfitting sul \emph{validation set} perché lo usa come ``oracolo'' per
scegliere l'architettura. Una rete che ha avuto fortuna in qualche
peculiarità del validation viene selezionata, ma quella fortuna non si
trasferisce al test set.
\end{concetto}

La penalità di complessità $\lambda$ è la difesa diretta: scoraggia il
GA dal selezionare reti grandi che hanno qualche frazione di punto di
accuracy in più ma molti parametri in più. Il Test 4 (sezione successiva)
mostra empiricamente questo effetto: senza penalità il GA seleziona reti
con migliaia di parametri; con $\lambda = 0.05$ seleziona la stessa
architettura ottima ma più snella.

% ─────────────────────────────────────────────────────────────
\subsection{Data leakage}
% ─────────────────────────────────────────────────────────────

Il \textbf{data leakage} è qualunque situazione in cui informazione del
validation o test set si ``infiltra'' nel processo di training,
gonfiando artificialmente le stime di performance. Le forme più comuni:

\begin{itemize}
    \item \textbf{Leakage da preprocessing} --- calcolare $\mu$, $\sigma$
          (StandardScaler) o altre statistiche sull'\emph{intero} dataset
          invece che solo sul training set. Le statistiche del test
          ``contaminano'' la trasformazione applicata al training. È
          oggetto del Test 9B.
    \item \textbf{Leakage temporale} --- in serie temporali, usare dati
          futuri per predire il passato. Non rilevante per Iris/Digits ma
          critico in finanza.
    \item \textbf{Bias ottimistico val vs test} --- selezionare un modello
          in base alla sua performance sul validation set e poi riportare
          quella stessa performance come ``risultato''. Il modello scelto è
          per costruzione il più fortunato sul validation, quindi la stima
          è ottimistica. Oggetto del Test 9A.
\end{itemize}

\begin{concetto}[la regola d'oro per evitare data leakage]
Tutto ciò che richiede statistiche dei dati (medie, varianze,
trasformazioni) si calcola \emph{esclusivamente} sul training set, e si
applica al validation/test \emph{senza ricalcolare nulla}. Il test set si
tocca \emph{una sola volta}, alla fine.
\end{concetto}

% ─────────────────────────────────────────────────────────────
\subsection{Esplorazione vs sfruttamento}
% ─────────────────────────────────────────────────────────────

Quasi ogni iperparametro del GA si può leggere come una scelta sul
trade-off tra esplorazione e sfruttamento:

\begin{center}
\begin{tabular}{p{3.4cm} p{4.6cm} p{4.6cm}}
    \toprule
    \rowcolor{headerblue}
    \textcolor{white}{\textbf{Iperparametro}} &
    \textcolor{white}{\textbf{Verso esplorazione}} &
    \textcolor{white}{\textbf{Verso sfruttamento}} \\
    \midrule
    population\_size  & grande (più diversità)        & piccolo (concentrazione) \\
    \rowcolor{rowgray}
    mutation\_rate    & alto (più variazioni)         & basso (preserva i buoni) \\
    tournament\_size  & piccolo (più chance ai meno forti) & grande (i migliori dominano) \\
    \rowcolor{rowgray}
    elitismo          & senza (più churn)             & con (preserva il top) \\
    \bottomrule
\end{tabular}
\end{center}

\begin{domanda}
\textit{``Cosa succede se il mutation rate è troppo alto?''}\\
Il GA si riduce a una ricerca quasi casuale. Ogni generazione modifica
così tanto i figli che le caratteristiche buone trovate dai genitori
vanno perse. Empiricamente, con $p_m \geq 0.4$ la \emph{mean accuracy}
della popolazione fatica a salire e diventa rumorosa: il segnale
evolutivo si perde.
\end{domanda}

% ─────────────────────────────────────────────────────────────
\subsection{Generalizzabilità: metodo vs risultato}
% ─────────────────────────────────────────────────────────────

\begin{domanda}
\textit{``L'architettura trovata su Digits funzionerebbe su altri dataset?''}\\
No. Il GA trova la migliore architettura \emph{per quel dataset specifico}.
L'architettura ottimale dipende dalla complessità del problema (numero di
feature, separabilità delle classi, quantità di dati). Il \emph{metodo}
(GA + NAS) è generalizzabile, ma il \emph{risultato} (l'architettura
trovata) è specifico per il dataset. Su problemi più complessi servirebbero
più generazioni, popolazione più grande e architetture più profonde. Il
Test 8 misura proprio questo, eseguendo lo stesso GA su quattro dataset
diversi e confrontando le architetture trovate.
\end{domanda}

% ─────────────────────────────────────────────────────────────
\subsection{Ottimizzazione annidata: due livelli}
% ─────────────────────────────────────────────────────────────

\begin{concetto}[GA + backpropagation = due ottimizzatori in cascata]
Nel progetto convivono due algoritmi di ottimizzazione che operano su
scale e parametri diversi:
\begin{itemize}
    \item \textbf{Livello esterno} (lento, su parametri \emph{discreti}):
          il GA ottimizza l'\emph{architettura} (numero di layer, neuroni,
          funzioni di attivazione). Lavora su una popolazione, una
          generazione alla volta.
    \item \textbf{Livello interno} (veloce, su parametri \emph{continui}):
          la backpropagation ottimizza i \emph{pesi} di ogni rete candidata.
          Lavora con SGD su un singolo individuo, una epoca alla volta.
\end{itemize}
La fitness del GA include la performance ottenuta dopo che il livello
interno è terminato. È un caso classico di \emph{bilevel optimization}.
\end{concetto}

\begin{domanda}
\textit{``Perché non far evolvere anche i pesi al GA?''}\\
Perché lo spazio dei pesi è continuo e ad altissima dimensionalità (decine
di migliaia di parametri reali per una rete media). Il GA è efficiente su
spazi discreti e di dimensione moderata, ma è enormemente inefficiente
sui pesi: la backpropagation, che usa l'informazione del gradiente, è
ordini di grandezza più rapida. La filosofia ``ognuno fa quello che sa
fare meglio'': GA sui parametri discreti (struttura), backpropagation sui
parametri continui (pesi).
\end{domanda}

\begin{riassunto}
\begin{itemize}
    \item Il GA fa overfitting sul \emph{validation set}, non sul training: la penalità $\lambda$ lo previene.
    \item Tre forme di data leakage da conoscere: preprocessing, temporale, bias val/test.
    \item Quasi ogni iperparametro del GA è un controllo del trade-off esplorazione/sfruttamento.
    \item Il metodo (GA + NAS) è generale, ma l'architettura trovata è specifica del dataset.
    \item Il progetto è un caso di \emph{ottimizzazione annidata}: GA sui parametri discreti, backpropagation sui pesi continui.
\end{itemize}
\end{riassunto}

\newpage

% =============================================================
\section{Esperimenti e Risultati}
% =============================================================

% ─────────────────────────────────────────────────────────────
\subsection{Metodologia: analisi di sensibilità}
% ─────────────────────────────────────────────────────────────

L'obiettivo degli esperimenti è rispondere a una domanda: \emph{quali
valori degli iperparametri del GA producono il miglior compromesso tra
qualità della soluzione, complessità della rete e costo computazionale?}

La metodologia adottata è l'\textbf{analisi di sensibilità}: si varia
\emph{un iperparametro alla volta} tenendo fissi tutti gli altri,
osservando l'effetto netto sulla \texttt{test\_accuracy} e su
\texttt{n\_params}. I test vengono eseguiti in un ordine preciso
(ottimizzazione sequenziale): ogni volta che un parametro viene fissato,
i test successivi usano quel valore aggiornato.

\subsubsection{Dataset principale: \fnt{load\_digits}}

Tutti i test di sensibilità sono condotti sul dataset \texttt{load\_digits}
di scikit-learn:
\begin{itemize}
    \item $1797$ immagini $8\times 8$ di cifre scritte a mano (riconoscibili
          come $0\dots 9$);
    \item $64$ feature continue (intensità dei pixel);
    \item $10$ classi (le cifre);
    \item complessità: alta. È stato scelto come dataset principale
          \emph{proprio perché} è abbastanza complesso da rendere
          significative le differenze fra architetture (su Iris, troppo
          semplice, ogni rete ragionevole arriva al $90\%$+).
\end{itemize}

Solo nel Test 8 si confrontano i risultati su quattro dataset diversi
(Iris, Wine, Breast Cancer, Digits) per testare la generalizzabilità del
metodo.

\subsubsection{Configurazione di base}

I valori di partenza sono volutamente conservativi (popolazione e
generazioni piccole per contenere il tempo di esecuzione):

\begin{center}
\begin{tabular}{ll}
    \toprule
    \rowcolor{headerblue}
    \textcolor{white}{\textbf{Parametro}} & \textcolor{white}{\textbf{Valore di partenza}} \\
    \midrule
    population\_size  & 20    \\
    \rowcolor{rowgray}
    generations       & 20    \\
    mutation\_rate    & 0.2   \\
    \rowcolor{rowgray}
    tournament\_size  & 5     \\
    epochs            & 150   \\
    \rowcolor{rowgray}
    learning\_rate    & 0.01  \\
    lambda\_          & 0.05  \\
    \rowcolor{rowgray}
    K                 & 5     \\
    seed              & 42    \\
    \bottomrule
\end{tabular}
\end{center}

L'ordine di ottimizzazione: \texttt{epochs} → \texttt{K} →
\texttt{learning\_rate} → \texttt{lambda} → \texttt{mutation\_rate} →
\texttt{population} → \texttt{generations}, seguito dai test 8
(generalizzabilità) e 9 (data leakage).

% ─────────────────────────────────────────────────────────────
\subsection{Test 1 — Numero di epoche}
% ─────────────────────────────────────────────────────────────

\textbf{Ipotesi.} Poche epoche $\Rightarrow$ la rete non converge
$\Rightarrow$ fitness rumorosa $\Rightarrow$ il GA seleziona architetture
``per caso''. Molte epoche $\Rightarrow$ stima accurata ma costosa.

\begin{center}
\begin{tabular}{lccr}
    \toprule
    \rowcolor{headerblue}
    \textcolor{white}{\textbf{epochs}} &
    \textcolor{white}{\textbf{test\_accuracy}} &
    \textcolor{white}{\textbf{tempo (s)}} &
    \textcolor{white}{\textbf{n\_params GA}} \\
    \midrule
    100   & 53.72\%  & 96.6   & 4\,735 \\
    \rowcolor{rowgray}
    250   & 78.33\%  & 96.8   & 4\,435 \\
    500   & 88.06\%  & 100.3  & 3\,910 \\
    \rowcolor{rowgray}
    1000  & 91.44\%  & 129.3  & 4\,360 \\
    \bottomrule
\end{tabular}
\end{center}

\textit{Baseline Digits: \textbf{682 parametri}.}

\textbf{Osservazioni.}
Con 100 epoche la \texttt{test\_accuracy} è $53.72\%$ — appena meglio di
un classificatore casuale su $10$ classi ($10\%$). La rete non ha avuto
tempo di convergere e il segnale di fitness è completamente inutilizzabile.
A 500 epoche si ottiene $88.06\%$ con un overhead di soli $4$ secondi
rispetto a 250: il salto è enorme. Da 500 a 1000 epoche il guadagno è di
$+3.4$ punti percentuali ma il tempo cresce del $29\%$ (rendimenti
decrescenti).

Con poche epoche il GA seleziona architetture con $4\,000$+ parametri,
molto più grandi della baseline ($682$): il segnale di fitness è rumoroso
e il GA non riesce a discriminare architetture compatte da quelle grandi,
quindi le architetture crescono liberamente. Solo con epoche sufficienti
la fitness diventa affidabile e la penalità $\lambda$ può fare effetto,
portando a reti più snelle.

\textbf{Scelta finale: \fnt{epochs = 500}.} Punto di massimo rendimento
marginale.

% ─────────────────────────────────────────────────────────────
\subsection{Test 2 — Numero di valutazioni $K$}
% ─────────────────────────────────────────────────────────────

\textbf{Ipotesi.} La fitness è la media di $K$ run indipendenti con seed
diversi. $K$ basso $\Rightarrow$ stima rumorosa. $K$ alto $\Rightarrow$
stima stabile ma ogni generazione costa $K$ volte di più.

\begin{center}
\begin{tabular}{lccr}
    \toprule
    \rowcolor{headerblue}
    \textcolor{white}{\textbf{K}} &
    \textcolor{white}{\textbf{mean\_test}} &
    \textcolor{white}{\textbf{std\_test}} &
    \textcolor{white}{\textbf{n\_params GA}} \\
    \midrule
    1   & 85.61\%   & 1.78\%  & 4\,135 \\
    \rowcolor{rowgray}
    3   & 85.70\%   & 0.93\%  & 3\,760 \\
    5   & 86.45\%   & 1.64\%  & 3\,910 \\
    \rowcolor{rowgray}
    10  & 86.86\%   & 0.79\%  & 3\,910 \\
    20  & 86.78\%   & 0.93\%  & 4\,735 \\
    \bottomrule
\end{tabular}
\end{center}

\textit{Baseline Digits: \textbf{682 parametri}.}

\textbf{Osservazioni.}
Con $K=1$ la deviazione standard sulla \texttt{test\_accuracy} (calcolata
su 5 seed diversi) è $1.78\%$: una rete può ricevere una valutazione
falsamente alta o bassa per pura casualità del seed. Con $K=3$ la
deviazione si dimezza a $0.93\%$ a costo solo triplicato. Da $K=3$ a
$K=10$ la \texttt{std} scende a $0.79\%$, miglioramento marginale di
$0.14$ punti che non giustifica il triplicarsi del costo.

In questo test il GA trova architetture con $3\,700$--$4\,700$ parametri,
molto più grandi della baseline ($682$): il learning rate è ancora a
$0.01$ (subottimale) e la fitness rumorosa impedisce alla penalità
$\lambda$ di operare efficacemente, lasciando crescere le architetture.

\textbf{Scelta finale: $\bm{K=3}$.} Dimezza il rumore di valutazione a
costo computazionale ragionevole.

% ─────────────────────────────────────────────────────────────
\subsection{Test 3 — Learning rate}
% ─────────────────────────────────────────────────────────────

\textbf{Ipotesi.} LR troppo basso $\Rightarrow$ la rete non converge in
500 epoche; LR troppo alto $\Rightarrow$ oscillazioni o instabilità
numerica.

\begin{center}
\begin{tabular}{llcr}
    \toprule
    \rowcolor{headerblue}
    \textcolor{white}{\textbf{learning\_rate}} &
    \textcolor{white}{\textbf{test\_accuracy}} &
    \textcolor{white}{\textbf{architettura trovata}} &
    \textcolor{white}{\textbf{n\_params GA}} \\
    \midrule
    0.001  & 38.15\%  & [(63, relu)]        & 4\,735 \\
    \rowcolor{rowgray}
    0.005  & 81.48\%  & [(64, relu)]        & 4\,810 \\
    0.01   & 87.31\%  & [(50, leaky\_relu)] & 3\,760 \\
    \rowcolor{rowgray}
    0.05   & 90.65\%  & [(24, relu)]        & 1\,810 \\
    0.1    & 90.74\%  & [(24, sigmoid)]     & 1\,810 \\
    \bottomrule
\end{tabular}
\end{center}

\textit{Baseline Digits: \textbf{682 parametri}.}

\textbf{Osservazioni.}
\begin{itemize}
    \item Con $\eta = 0.001$ la rete non converge in $500$ epoche
          ($38.15\%$ è poco più di un classificatore casuale): i passi di
          gradient descent sono troppo piccoli.
    \item Con $\eta = 0.05$ e $\eta = 0.1$ si ottiene accuracy quasi
          identica, ma con $\eta=0.1$ il GA seleziona un'architettura con
          \fnt{sigmoid}: questa funzione satura facilmente quando la LR è
          alta, introducendo instabilità nei test successivi.
    \item Con $\eta = 0.05$ il GA trova \fnt{[(24, relu)]} in modo stabile
          e riproducibile.
    \item \textbf{Il learning rate influenza i parametri trovati:} con LR
          basso ($\le 0.01$) il GA seleziona architetture con $3\,700$--
          $4\,800$ parametri (fino a $7\times$ la baseline). Con LR
          ottimale ($0.05$) converge a $1\,810$ parametri. La ragione è
          che con LR basso il segnale di fitness è rumoroso e il GA non
          riesce a discriminare reti compatte da quelle grandi --- le
          architetture crescono senza controllo.
\end{itemize}

\begin{concetto}[il learning rate è il parametro più critico]
Su tutti i parametri testati, il learning rate produce la differenza più
ampia tra caso peggiore e caso migliore (oltre $50$ punti percentuali!).
È il parametro su cui si gioca tutto: una scelta sbagliata rende il GA
completamente inutile, perché il segnale di fitness diventa rumore puro.
\end{concetto}

\textbf{Scelta finale: \fnt{learning\_rate = 0.05}.} Massima accuracy
stabile con architettura più snella.

% ─────────────────────────────────────────────────────────────
\subsection{Test 4 — Penalità di complessità $\lambda$}
% ─────────────────────────────────────────────────────────────

\textbf{Ipotesi.} $\lambda = 0$ $\Rightarrow$ il GA seleziona reti grandi
(overfitting sul validation, gap val/test elevato). $\lambda$ alto
$\Rightarrow$ reti minuscole con accuracy insufficiente.

\begin{center}
\begin{tabular}{lccl}
    \toprule
    \rowcolor{headerblue}
    \textcolor{white}{\textbf{$\lambda$}} &
    \textcolor{white}{\textbf{test\_accuracy}} &
    \textcolor{white}{\textbf{n\_params}} &
    \textcolor{white}{\textbf{architettura}} \\
    \midrule
    0.0  & 90.65\%  & 1\,810   & [(24, relu)]            \\
    \rowcolor{rowgray}
    0.01 & 90.65\%  & 1\,810   & [(24, relu)]            \\
    0.05 & 90.65\%  & 1\,810   & [(24, relu)]            \\
    \rowcolor{rowgray}
    0.1  & 88.80\%  & 1\,135   & [(15, leaky\_relu)]     \\
    0.2  & 86.20\%  &   760    & [(10, relu)]            \\
    \bottomrule
\end{tabular}
\end{center}

\textit{Baseline Digits: \textbf{682 parametri}.}

\textbf{Osservazioni.}
\begin{itemize}
    \item Da $\lambda = 0$ a $\lambda = 0.05$ la \texttt{test\_accuracy}
          è identica e la stessa architettura viene selezionata ($1\,810$
          parametri). La penalità c'è ma non è abbastanza forte da
          distorcere la scelta: $1\,810$ parametri sono più che necessari
          per Digits, ma la penalità logaritmica non li penalizza
          abbastanza da spingersi verso reti più piccole.
    \item Da $\lambda = 0.1$ in poi la \texttt{test\_accuracy} cala
          ($-1.85\%$, poi $-4.45\%$) perché il GA preferisce reti troppo
          piccole: $1\,135$ e $760$ parametri sono insufficienti per la
          complessità di Digits.
    \item La baseline con $682$ parametri ottiene solo $58.33\%$, a
          conferma che su Digits servono più parametri di quelli che la
          baseline fissa può offrire.
\end{itemize}

\textbf{Scelta finale: $\bm{\lambda = 0.05}$.} Mantiene tutta l'accuracy
disponibile e introduce un incentivo alla semplicità che agisce sui casi
borderline (es. quando due architetture hanno accuracy molto simile,
quella con meno parametri viene preferita).

\begin{domanda}
\textit{``Perché penalità logaritmica e non lineare?''}\\
Con la penalità lineare $\lambda \cdot N_{\text{params}}$ il valore
ottimale di $\lambda$ dipende dalla scala assoluta delle reti generate.
Una stessa $\lambda$ su una rete da 100 parametri quasi non fa nulla,
sulla stessa rete a 10000 parametri annienta tutto. Con $\log_{10}$ il
fattore di crescita della penalità è mite (passare da 100 a 10000
parametri aumenta la penalità solo di $2$), rendendo $\lambda$ stabile
indipendentemente dalla scala.
\end{domanda}

% ─────────────────────────────────────────────────────────────
\subsection{Test 5 — Mutation rate}
% ─────────────────────────────────────────────────────────────

\textbf{Ipotesi.} Rate basso $\Rightarrow$ il GA si appoggia solo al
crossover, rischia ottimi locali. Rate alto $\Rightarrow$ il GA si riduce
a ricerca casuale.

\begin{center}
\begin{tabular}{lccl}
    \toprule
    \rowcolor{headerblue}
    \textcolor{white}{\textbf{$p_m$}} &
    \textcolor{white}{\textbf{test\_accuracy}} &
    \textcolor{white}{\textbf{n\_params}} &
    \textcolor{white}{\textbf{architettura}} \\
    \midrule
    0.0  & 89.07\%  & 2\,110 & [(28, leaky\_relu)] \\
    \rowcolor{rowgray}
    0.05 & 89.07\%  & 2\,110 & [(28, leaky\_relu)] \\
    0.1  & 90.19\%  & 1\,885 & [(25, relu)]        \\
    \rowcolor{rowgray}
    0.2  & 90.65\%  & 1\,810 & [(24, relu)]        \\
    0.4  & 90.19\%  & 1\,885 & [(25, relu)]        \\
    \rowcolor{rowgray}
    0.8  & 91.11\%  & 1\,735 & [(23, relu)]        \\
    \bottomrule
\end{tabular}
\end{center}

\textit{Baseline Digits: \textbf{682 parametri}.}

\textbf{Osservazioni.}
\begin{itemize}
    \item Con $p_m = 0$ e $p_m = 0.05$ il GA converge sulla stessa
          architettura ($28$ neuroni, leaky\_relu): senza mutazione
          significativa, il crossover da solo rimane intrappolato in un
          ottimo locale.
    \item $p_m = 0.2$ è il picco tra i valori \emph{stabili}.
    \item $p_m = 0.8$ produce una \texttt{test\_accuracy} apparentemente
          migliore, ma con l'$80\%$ dei geni mutati ad ogni passo il GA
          non costruisce più sull'evoluzione precedente: i risultati
          sono altamente dipendenti dal seed e poco riproducibili.
\end{itemize}

\textbf{Scelta finale: $\bm{p_m = 0.2}$.} Il rate più conservativo che
raggiunge il picco di \texttt{test\_accuracy} con comportamento stabile.

% ─────────────────────────────────────────────────────────────
\subsection{Test 6 — Dimensione della popolazione}
% ─────────────────────────────────────────────────────────────

\textbf{Ipotesi.} Popolazione piccola $\Rightarrow$ ricerca ristretta,
convergenza prematura. Popolazione grande $\Rightarrow$ esplorazione
ampia ma con generazioni fissate ogni individuo viene raffinato meno.

\begin{center}
\begin{tabular}{lccc}
    \toprule
    \rowcolor{headerblue}
    \textcolor{white}{\textbf{population}} &
    \textcolor{white}{\textbf{test\_accuracy}} &
    \textcolor{white}{\textbf{gap val$-$test}} &
    \textcolor{white}{\textbf{n\_params}} \\
    \midrule
    5    & 87.69\%  & 3.33\%  & 1\,585 \\
    \rowcolor{rowgray}
    10   & 88.98\%  & 1.85\%  & 1\,510 \\
    20   & 90.65\%  & 0.83\%  & 1\,810 \\
    \rowcolor{rowgray}
    30   & 89.07\%  & 2.41\%  & 2\,110 \\
    50   & 89.63\%  & 1.30\%  & 1\,360 \\
    \rowcolor{rowgray}
    70   & 90.00\%  & 1.67\%  & 1\,960 \\
    100  & 88.80\%  & 0.94\%  & 1\,135 \\
    \bottomrule
\end{tabular}
\end{center}

\textit{Baseline Digits: \textbf{682 parametri}.}

\textbf{Osservazioni.}
La relazione \emph{non} è monotona: c'è un picco netto a popolazione
$=20$, sia in \texttt{test\_accuracy} che in gap val/test minimo. Il gap
val/test minimo a $0.83\%$ è una doppia conferma --- non è un artefatto
della validation set.

Con popolazioni grandi e numero di generazioni fisso, ogni generazione
esplora più candidati ma li raffina meno (la pressione selettiva è
diluita). Con popolazioni piccole, manca diversità per uscire da ottimi
locali. Il punto di equilibrio è $20$.

\textbf{Scelta finale: \fnt{population = 20}.}

% ─────────────────────────────────────────────────────────────
\subsection{Test 7 — Numero di generazioni}
% ─────────────────────────────────────────────────────────────

\textbf{Ipotesi.} Poche generazioni $\Rightarrow$ il GA non ha avuto
tempo di convergere. Troppe $\Rightarrow$ costo computazionale inutile
perché la fitness è già stabile.

\begin{center}
\begin{tabular}{lclr}
    \toprule
    \rowcolor{headerblue}
    \textcolor{white}{\textbf{generations}} &
    \textcolor{white}{\textbf{test\_accuracy}} &
    \textcolor{white}{\textbf{architettura}} &
    \textcolor{white}{\textbf{n\_params GA}} \\
    \midrule
    5    & 88.61\%  & [(52, relu)] & 3\,910 \\
    \rowcolor{rowgray}
    10   & 90.65\%  & [(24, relu)] & 1\,810 \\
    20   & 90.65\%  & [(24, relu)] & 1\,810 \\
    \rowcolor{rowgray}
    30   & 90.65\%  & [(24, relu)] & 1\,810 \\
    50   & 90.65\%  & [(24, relu)] & 1\,810 \\
    \rowcolor{rowgray}
    70   & 90.65\%  & [(24, relu)] & 1\,810 \\
    100  & 90.65\%  & [(24, relu)] & 1\,810 \\
    \bottomrule
\end{tabular}
\end{center}

\textit{Baseline Digits: \textbf{682 parametri}.}

\textbf{Osservazioni.}
La convergenza è sorprendentemente rapida: già a $\text{gen} = 10$ il GA
trova la stessa architettura ottimale che trova anche a $\text{gen} = 100$.
Il salto da $5$ a $10$ è l'unico significativo ($+2.04\%$); da $10$ in poi
il risultato è identico. Da $\text{gen}=10$ in poi anche il numero di
parametri si stabilizza a $1\,810$: l'architettura è convergita.

\textbf{Scelta finale: \fnt{generations = 50}.} Ampio margine di sicurezza
sul punto di convergenza reale ($\text{gen}=10$), senza sprecare tempo
oltre.

% ─────────────────────────────────────────────────────────────
\subsection{Test 8 — Generalizzabilità su dataset diversi}
% ─────────────────────────────────────────────────────────────

\textbf{Ipotesi.} La configurazione ottimale trovata su Digits non è
sovra-specializzata per quel dataset, ma funziona ragionevolmente anche
su altri.

\begin{center}
\begin{tabular}{lccccl}
    \toprule
    \rowcolor{headerblue}
    \textcolor{white}{\textbf{dataset}} &
    \textcolor{white}{\textbf{feat.}} &
    \textcolor{white}{\textbf{cl.}} &
    \textcolor{white}{\textbf{GA test\%}} &
    \textcolor{white}{\textbf{Baseline\%}} &
    \textcolor{white}{\textbf{vince (acc.)?}} \\
    \midrule
    iris           &  4 &  3 & 92.22\%   & 94.44\%   & no         \\
    \rowcolor{rowgray}
    wine           & 13 &  3 & 100.00\%  & 100.00\%  & pari$^*$   \\
    breast\_cancer & 30 &  2 & 95.61\%   & 96.20\%   & no         \\
    \rowcolor{rowgray}
    digits         & 64 & 10 & 90.65\%   & 58.33\%   & sì (+32\%) \\
    \bottomrule
\end{tabular}
\end{center}

$^*$\textit{Accuracy identica ma architetture molto diverse per numero di parametri: vedere tabella sotto.}

\medskip
\textbf{Confronto parametri: GA vs baseline.}

\begin{center}
\begin{tabular}{lrrlrr}
    \toprule
    \rowcolor{headerblue}
    \textcolor{white}{\textbf{dataset}} &
    \textcolor{white}{\textbf{params GA}} &
    \textcolor{white}{\textbf{architettura GA}} &
    \textcolor{white}{\textbf{}} &
    \textcolor{white}{\textbf{params BL}} &
    \textcolor{white}{\textbf{rapporto}} \\
    \midrule
    iris           &  43  & [(5, leaky\_relu)]  & & 139  & GA usa $31\%$  \\
    \rowcolor{rowgray}
    wine           & 122  & [(7, relu)]         & & 211  & GA usa $58\%$  \\
    breast\_cancer & 167  & [(5, relu)]         & & 338  & GA usa $49\%$  \\
    \rowcolor{rowgray}
    digits         & 1\,810 & [(24, relu)]      & & 682  & GA usa $266\%$ \\
    \bottomrule
\end{tabular}
\end{center}

\textbf{Osservazioni.}
\begin{itemize}
    \item Su \emph{Iris} e \emph{Breast Cancer} la baseline batte il GA
          di pochi decimi di punto in accuracy, ma usa $3$--$4.5\times$
          più parametri. Il GA trova soluzioni più parsimonious (43 vs 139
          e 167 vs 338): a parità quasi completa di performance, la rete
          più piccola è da preferire perché generalizza meglio su dati
          rumorosi.
    \item Su \emph{Wine} entrambi raggiungono il $100\%$, ma il GA ottiene
          lo stesso risultato con soli $122$ parametri contro i $211$ della
          baseline ($-42\%$). La parità di accuracy \emph{non significa}
          parità di soluzioni: la baseline ha $89$ parametri in eccesso
          rispetto a quanti ne servono realmente. Questi parametri ``in
          più'' non portano benefici, e su dati più rumorosi o con meno
          campioni aumenterebbero il rischio di overfitting (vedi
          \S\,\ref{ssec:overfitting}).
    \item Su \emph{Digits} il GA batte la baseline di $32$ punti
          percentuali e giustamente usa $1\,810$ parametri ($2.7\times$
          la baseline). La baseline con $682$ parametri è
          \emph{sottodimensionata} per la complessità del problema
          ($64$ feature, $10$ classi): non ha abbastanza capacità e
          rimane a $58.33\%$. Il GA esplora lo spazio e trova la
          dimensione giusta.
\end{itemize}

\begin{concetto}[la NAS aggiunge valore quando il problema è non triviale]
Su dataset semplici, una rete fissa ragionevole è già adeguata: la NAS
non guadagna molto in accuracy, ma trova comunque reti \emph{più
compatte}. Su dataset complessi, la ricerca automatica può fare una
differenza enorme sia in accuracy che nel trovare la capacità giusta.
Questo è anche il motivo per cui la NAS è particolarmente utile in
applicazioni industriali (visione artificiale, NLP) dove i problemi sono
complessi e la dimensione ottimale della rete non è ovvia.
\end{concetto}

\textbf{Osservazione di rilievo.} Su \emph{tutti} i dataset il GA trova
architetture a un solo hidden layer. Non è un limite del GA: è una
risposta ai dati. Per questi problemi un singolo layer nascosto è
sufficiente, e la penalità $\lambda$ disincentiva la profondità inutile.

% ─────────────────────────────────────────────────────────────
\subsection{Test 9A — Bias ottimistico val vs test}
% ─────────────────────────────────────────────────────────────

\textbf{Ipotesi.} Il GA seleziona l'architettura migliore in base
all'accuracy sul \emph{validation} set. Riportare quella stessa accuracy
come ``risultato'' del GA è ottimistico per costruzione.

\begin{center}
\begin{tabular}{lc}
    \toprule
    \rowcolor{headerblue}
    \textcolor{white}{\textbf{metrica}} &
    \textcolor{white}{\textbf{valore (media su 10 seed)}} \\
    \midrule
    \texttt{val\_accuracy}     & 91.33\%  \\
    \rowcolor{rowgray}
    \texttt{test\_accuracy}    & 87.01\%  \\
    \textbf{gap medio}         & \textbf{4.33\%} \\
    \bottomrule
\end{tabular}
\end{center}

\textbf{Osservazioni.}
Il gap di $4.33\%$ è sistematico. Su 10 seed diversi varia da $1.39\%$ a
$9.35\%$, ma è sempre positivo (la val è sempre più alta della test).
Conclusione: la \texttt{test\_accuracy} è la stima onesta da riportare;
la \texttt{best\_accuracy} (val) è ottimistica per costruzione.

\begin{concetto}[lezione operativa]
Quando si fa NAS o tuning di iperparametri tramite validation set, è
\emph{obbligatorio} avere un test set separato per la stima finale.
Riportare la metrica usata per scegliere significa riportare il caso
fortunato, non la performance attesa.
\end{concetto}

% ─────────────────────────────────────────────────────────────
\subsection{Test 9B — Leakage da preprocessing}
% ─────────────────────────────────────────────────────────────

\textbf{Ipotesi.} Calcolare le statistiche dello \fnt{StandardScaler}
sull'intero dataset (train + val + test) anziché solo sul training
introduce leakage e gonfia la stima.

\begin{center}
\begin{tabular}{lc}
    \toprule
    \rowcolor{headerblue}
    \textcolor{white}{\textbf{scenario}} &
    \textcolor{white}{\textbf{test\_accuracy}} \\
    \midrule
    scaler corretto (solo train)   & 75.33\% \\
    \rowcolor{rowgray}
    scaler su tutto il dataset     & 74.00\% \\
    \bottomrule
\end{tabular}
\end{center}

\textbf{Osservazioni.}
In questo caso specifico, il leakage non ha gonfiato i risultati (anzi
li ha leggermente peggiorati). Il motivo è che il dataset usato per
questo test (Iris) è piccolo e omogeneo: la distribuzione dei dati è
molto simile in ciascuno split, quindi le statistiche calcolate
sull'intero dataset coincidono praticamente con quelle calcolate solo sul
training.

Su dataset più grandi o con distribuzione eterogenea il leakage sarebbe
più evidente: l'esperimento documenta il fatto che la \emph{prassi}
(scaler solo sul training) va comunque adottata sempre, anche quando
nei casi piccoli non sembra fare differenza.

% ─────────────────────────────────────────────────────────────
\subsection{Riepilogo della configurazione ottimale}
% ─────────────────────────────────────────────────────────────

\begin{center}
\begin{tabular}{ll>{\raggedright\arraybackslash}p{6.5cm}}
    \toprule
    \rowcolor{headerblue}
    \textcolor{white}{\textbf{Parametro}} &
    \textcolor{white}{\textbf{Valore}} &
    \textcolor{white}{\textbf{Motivazione}} \\
    \midrule
    epochs          & 500    & Minimo per segnale fitness affidabile (rendimenti decrescenti oltre)  \\
    \rowcolor{rowgray}
    K               & 3      & Dimezza la varianza di K=1 a costo minimo                            \\
    learning\_rate  & 0.05   & Picco di test\_accuracy con convergenza stabile (più critico di tutti) \\
    \rowcolor{rowgray}
    lambda          & 0.05   & Penalizza la complessità senza ridurre l'accuracy                     \\
    mutation\_rate  & 0.2    & Picco stabile, evita ottimi locali senza diventare ricerca casuale    \\
    \rowcolor{rowgray}
    population      & 20     & Picco di test\_accuracy + minimo gap val/test                         \\
    generations     & 50     & Ampio margine sul punto di convergenza reale (gen=10)                 \\
    \bottomrule
\end{tabular}
\end{center}

\textbf{Risultato finale su Digits:} architettura \texttt{[(24, relu)]} —
1 hidden layer, $24$ neuroni, ReLU. \\
\textbf{Test accuracy:} $90.65\%$ \;vs\; baseline $58.33\%$
($+32.32$ punti percentuali).

\begin{riassunto}
\begin{itemize}
    \item Il \emph{learning rate} è di gran lunga il parametro più critico: una scelta sbagliata rende il GA inutile.
    \item Il numero di \emph{epoche} controlla la qualità del segnale di fitness; troppo poche $\Rightarrow$ rumore puro.
    \item $K$ (numero di valutazioni) e $\lambda$ (penalità) controllano rispettivamente rumore e complessità.
    \item \emph{population} e \emph{generations} contano meno di quanto ci si aspetterebbe: il GA converge presto.
    \item Il GA aggiunge valore reale solo su dataset \emph{complessi}: su Iris la baseline è già adeguata.
    \item Test set indispensabile per stima onesta: la val accuracy è ottimistica di $\sim 4\%$ in media.
\end{itemize}
\end{riassunto}

\newpage

% =============================================================
\section{Discussione e Conclusioni}
% =============================================================

% ─────────────────────────────────────────────────────────────
\subsection{Cosa il progetto dimostra}
% ─────────────────────────────────────────────────────────────

Il progetto dimostra empiricamente, su un dataset non triviale, che
\textbf{un Algoritmo Genetico è uno strumento efficace per la Neural
Architecture Search}. Sul dataset Digits il GA scopre automaticamente
un'architettura ($[(24, \text{relu})]$) che batte una baseline manuale
ragionevole di $32$ punti percentuali. La ricerca converge in poche
generazioni e produce architetture compatte (un solo layer nascosto), un
risultato non ovvio a priori.

Più in generale, il progetto mostra che:
\begin{itemize}
    \item È possibile combinare due paradigmi diversi dell'AI (ricerca
          evolutiva ed apprendimento supervisionato) in un sistema
          coerente.
    \item L'\emph{ottimizzazione annidata} (GA sopra, backpropagation
          sotto) funziona purché ognuno dei due livelli sia configurato
          correttamente.
    \item Il principio di \emph{regolarizzazione} via penalità di
          complessità mantiene le architetture trovate snelle e
          robuste, riducendo l'overfitting sul validation set.
\end{itemize}

% ─────────────────────────────────────────────────────────────
\subsection{Limiti del progetto}
% ─────────────────────────────────────────────────────────────

\begin{itemize}
    \item \textbf{Generalizzabilità del risultato.} Le architetture
          ottimali trovate sono \emph{specifiche per il dataset} su cui
          sono state cercate. Il \emph{metodo} è generale; il
          \emph{risultato} no.
    \item \textbf{Spazio di ricerca limitato.} Il GA cerca solo nel
          sottospazio delle reti feed-forward MLP. Architetture
          convoluzionali, ricorrenti o con \emph{skip connections} non
          sono nello spazio di ricerca. Questo è una limitazione di
          implementazione, non concettuale.
    \item \textbf{Iperparametri di addestramento esclusi.} Learning rate
          ed epochs sono iperparametri della backpropagation, fissati
          all'esterno del GA. Una versione più ambiziosa potrebbe
          includerli nel cromosoma.
    \item \textbf{Costo computazionale.} L'unica vera ``commodità'' del
          progetto. Su un singolo dataset con la configurazione finale, una
          run completa richiede dell'ordine dei minuti su CPU moderna.
          Per dataset molto grandi (ImageNet) approcci più sofisticati
          (NAS basato su reinforcement learning, weight sharing) sono
          necessari.
    \item \textbf{Variabilità tra run.} Anche con $K=3$, esecuzioni con
          seed diversi possono trovare architetture leggermente diverse.
          Il sistema è stocastico per natura: non c'è un'unica ``risposta
          giusta'' a ogni esecuzione.
\end{itemize}

% ─────────────────────────────────────────────────────────────
\subsection{Sviluppi futuri}
% ─────────────────────────────────────────────────────────────

\begin{itemize}
    \item \textbf{Estendere il cromosoma.} Includere learning rate,
          dimensione del batch, dropout rate come geni. Il GA
          ottimizzerebbe simultaneamente architettura e training.
    \item \textbf{Operatori di crossover più intelligenti.} Crossover
          consapevoli della semantica (es. scambiare layer simili in
          posizione simile) potrebbero accelerare la convergenza.
    \item \textbf{Funzione di fitness multi-obiettivo.} Invece di
          combinare accuracy e complessità in un'unica fitness scalare,
          esplorare il fronte di Pareto con un GA multi-obiettivo
          (NSGA-II).
    \item \textbf{Fitness stimata.} Il collo di bottiglia è il training
          completo per ogni valutazione. Un \emph{fitness predictor} (a
          sua volta una rete neurale) potrebbe stimare la fitness senza
          training, accelerando l'evoluzione di ordini di grandezza.
    \item \textbf{Spazi di ricerca più ricchi.} Estendere il GA a gestire
          architetture non lineari (skip connection, branching).
\end{itemize}

\newpage

% =============================================================
\section{Domande tipo orale (con risposte ragionate)}
% =============================================================

Questa sezione raccoglie le domande più probabili per l'orale, con
risposte sintetiche ma complete. Le domande sono organizzate per
argomento.

% ─────────────────────────────────────────────────────────────
\subsection{Sul progetto in generale}
% ─────────────────────────────────────────────────────────────

\begin{domanda}
\textit{``Spiegami il tuo progetto in due minuti.''}\\
Ho implementato un sistema di Neural Architecture Search: un Algoritmo
Genetico cerca automaticamente la migliore architettura di una rete
neurale (numero di layer, neuroni, funzioni di attivazione) per un dato
dataset. Ogni architettura candidata viene valutata addestrandola con
backpropagation e misurando l'accuratezza sul validation set, penalizzata
per la complessità della rete. Sul dataset Digits il sistema ha trovato
un'architettura con un singolo layer da $24$ neuroni ReLU, $90.65\%$ di
accuracy sul test set, contro il $58.33\%$ di una baseline manuale. Sia la
rete che il GA sono scritti da zero in numpy.
\end{domanda}

\begin{domanda}
\textit{``Perché hai scelto un GA invece di altri metodi di NAS?''}\\
Per tre motivi: (1) il GA è naturale per spazi di ricerca discreti come
quello delle architetture; (2) gestisce nativamente cromosomi di
lunghezza variabile (architetture con numero di layer diverso); (3) è
algoritmicamente semplice da implementare da zero, coerente con
l'obiettivo didattico del progetto. Alternative come grid search sono
esponenzialmente costose, e reinforcement learning per NAS richiede
infrastruttura molto più pesante.
\end{domanda}

\begin{domanda}
\textit{``Quali sono le componenti principali del codice?''}\\
Due moduli: \texttt{neural\_network} contiene la classe della rete con
forward, backward e una libreria di funzioni di attivazione;
\texttt{genetic\_algorithm} contiene la classe del GA con i metodi
privati per fitness, selezione, crossover, mutazione, e il metodo
pubblico \texttt{run()} che orchestra il ciclo evolutivo. C'è poi un
\texttt{main.py} che fa preprocessing del dataset, esegue il GA e produce
i grafici, e una serie di script di test che variano un parametro alla
volta per l'analisi di sensibilità.
\end{domanda}

% ─────────────────────────────────────────────────────────────
\subsection{Su reti neurali}
% ─────────────────────────────────────────────────────────────

\begin{domanda}
\textit{``Cos'è la backpropagation, in poche frasi?''}\\
È l'algoritmo che calcola il gradiente della loss rispetto a tutti i
pesi della rete, applicando la regola della catena layer per layer
all'indietro a partire dall'output. Una volta ottenuto il gradiente, i
pesi vengono aggiornati nella direzione opposta tramite gradient descent.
Si chiama \emph{back}-propagation perché il calcolo procede
dall'\emph{output verso l'input}, riusando i risultati intermedi
salvati durante il forward pass per essere efficiente.
\end{domanda}

\begin{domanda}
\textit{``A cosa serve la funzione di attivazione non lineare?''}\\
A rendere la rete capace di rappresentare funzioni non lineari. Senza
attivazioni non lineari, una rete con $N$ layer è matematicamente
equivalente a una rete con un singolo layer lineare: la composizione di
trasformazioni lineari è lineare. È la non linearità a dare alla rete il
suo potere rappresentativo (universal approximation theorem).
\end{domanda}

\begin{domanda}
\textit{``Perché si usa softmax in output e non sigmoid?''}\\
La softmax produce una distribuzione di probabilità sulle classi (i
valori sommano a $1$): è il modo naturale di esprimere ``il modello
crede che questo input appartenga alla classe $i$ con probabilità $p_i$''.
Sigmoid produrrebbe valori indipendenti per ogni classe, non collegati
fra loro. Per classificazione binaria si usa sigmoid; per multiclasse
softmax.
\end{domanda}

\begin{domanda}
\textit{``Cos'è il vanishing gradient e come si previene?''}\\
È il fenomeno per cui i gradienti diventano sempre più piccoli scendendo
verso i layer iniziali della rete, fino a impedirne l'apprendimento.
Cause: funzioni di attivazione che saturano (sigmoid, tanh) e cattiva
inizializzazione dei pesi. Si previene usando ReLU/Leaky ReLU (che non
saturano per input positivi), inizializzazione di He o Xavier (che
mantiene la varianza dei segnali costante attraverso i layer), e
gradient clipping (protegge dal caso opposto, exploding gradient).
\end{domanda}

% ─────────────────────────────────────────────────────────────
\subsection{Su algoritmi genetici}
% ─────────────────────────────────────────────────────────────

\begin{domanda}
\textit{``Spiegami il ciclo di un algoritmo genetico.''}\\
Si parte da una popolazione casuale di soluzioni candidate. A ogni
generazione: (1) si calcola la fitness di ogni individuo; (2) si
selezionano i genitori (tipicamente con torneo); (3) i genitori si
combinano via crossover producendo dei figli; (4) i figli vengono
mutati con probabilità $p_m$; (5) la nuova popolazione sostituisce la
vecchia, eventualmente preservando il migliore con elitismo. Si itera
finché la fitness converge o si raggiunge un numero massimo di
generazioni.
\end{domanda}

\begin{domanda}
\textit{``Cos'è l'elitismo e perché è utile?''}\\
È la tecnica di copiare il migliore individuo della generazione corrente
nella nuova senza farlo passare per crossover/mutazione. Garantisce che
la fitness massima non possa peggiorare da una generazione all'altra:
quello che si è trovato non si perde più. Il costo è una piccola
riduzione di diversità (uno slot in meno per nuovi individui), ma con
popolazioni di $20$+ è trascurabile.
\end{domanda}

\begin{domanda}
\textit{``Esplorazione vs sfruttamento --- come si bilanciano nel tuo GA?''}\\
Diversi parametri controllano questo trade-off. Il \texttt{mutation\_rate}
($0.2$) inietta abbastanza casualità da evitare ottimi locali ma non
abbastanza da degradare il GA in ricerca casuale. Il \texttt{tournament\_size}
($5$) è abbastanza piccolo da dare chance anche agli individui mediocri
(esplorazione) ma abbastanza grande da far emergere i migliori
(sfruttamento). L'\emph{elitismo} aggiunge sfruttamento puro. La
\emph{popolazione} di $20$ individui è abbastanza ampia da mantenere
diversità senza diluire la pressione selettiva.
\end{domanda}

\begin{domanda}
\textit{``Cosa succede se aumenti molto il tournament size?''}\\
Aumenta la pressione selettiva: i migliori vincono quasi sempre i
tornei, gli individui mediocri quasi mai. Il GA converge più rapidamente
ma su un ottimo (anche locale): la popolazione perde diversità in fretta
e il GA smette di esplorare. È l'opposto della ricerca casuale --- una
``hill climbing su steroidi''.
\end{domanda}

% ─────────────────────────────────────────────────────────────
\subsection{Su NAS e fitness}
% ─────────────────────────────────────────────────────────────

\begin{domanda}
\textit{``Perché K = 3 valutazioni per ogni individuo?''}\\
Perché i pesi della rete sono inizializzati casualmente a ogni
valutazione, e con SGD a campioni casuali il training stesso è
stocastico. Conseguenza: la stessa architettura, valutata due volte di
seguito, può dare accuracy diverse di $1$--$2\%$ per pura fortuna.
Mediare su $K=3$ run dimezza la deviazione standard a costo
ragionevole. Aumentare $K$ riduce ulteriormente la varianza ma con
rendimenti decrescenti rapidi.
\end{domanda}

\begin{domanda}
\textit{``Perché penalizzi la complessità della rete? Cosa succederebbe senza?''}\\
Senza penalità il GA tende a far crescere le architetture, perché reti
più grandi hanno tipicamente accuracy leggermente superiore sul
validation set (più capacità di adattamento). Questo causa due problemi:
(1) overfitting sul validation, con gap val/test elevato; (2) reti
inutilmente costose. La penalità $\lambda \cdot \log_{10}(N_{\text{params}})$
agisce come regolarizzatore: scoraggia la complessità non giustificata
da un guadagno sostanziale in accuracy.
\end{domanda}

\begin{domanda}
\textit{``Perché la penalità è logaritmica e non lineare?''}\\
Con la penalità lineare $\lambda \cdot N_{\text{params}}$ il valore
ottimale di $\lambda$ dipende fortemente dalla scala delle reti generate
in quella esecuzione: una stessa $\lambda$ non penalizza praticamente
nulla una rete piccola e annienta una rete grande. Con $\log_{10}$ la
penalità cresce in modo molto più mite (passare da 100 a 10000
parametri aumenta il termine solo di $2$), rendendo $\lambda$ stabile
indipendentemente dalla scala.
\end{domanda}

\begin{domanda}
\textit{``Come fai a confrontare reti con architetture diverse?''}\\
Tramite la fitness. Ogni rete viene valutata con la stessa procedura:
costruita con i suoi parametri di architettura, addestrata con
backpropagation per $E$ epoche, valutata in accuracy sul validation set.
La fitness è poi accuracy meno la penalità $\lambda \cdot
\log_{10}(N_{\text{params}})$, che permette di confrontare reti di
dimensioni diverse: una rete piccola con accuracy del $90\%$ può avere
fitness più alta di una rete grande con accuracy del $92\%$.
\end{domanda}

% ─────────────────────────────────────────────────────────────
\subsection{Su validazione e leakage}
% ─────────────────────────────────────────────────────────────

\begin{domanda}
\textit{``Differenza tra validation set e test set?''}\\
Il \emph{validation} si usa \emph{durante} la ricerca, per scegliere fra
configurazioni alternative. Il \emph{test} si usa \emph{una sola volta}
alla fine, per stimare onestamente le prestazioni. Se si usa il test set
durante la ricerca, si introduce data leakage e la stima diventa
ottimistica. Nel mio progetto: il GA misura la fitness sul validation,
e solo l'architettura vincitrice viene rivalutata sul test set, da cui
la metrica finale.
\end{domanda}

\begin{domanda}
\textit{``Cos'è il data leakage e come lo eviti?''}\\
È qualunque situazione in cui informazione del validation/test ``filtra''
nel processo di training, gonfiando le stime di performance. Le forme
più tipiche sono: (1) calcolare statistiche di preprocessing (es.
StandardScaler) sull'intero dataset invece che solo sul training; (2)
selezionare un modello con il validation e poi riportare la sua accuracy
sul validation come risultato (bias ottimistico). Lo evito calcolando
ogni statistica solo sul training, applicandola identica a val/test, e
riportando come risultato finale solo la \texttt{test\_accuracy}.
\end{domanda}

\begin{domanda}
\textit{``Cosa hai trovato col Test 9A sul bias val/test?''}\\
Su $10$ seed diversi, la val accuracy media è del $91.33\%$ ma la test
accuracy media è solo $87.01\%$: un gap sistematico di $4.33\%$. Questo
conferma che la metrica usata per scegliere (validation) non è una
stima onesta delle prestazioni: è ottimistica per costruzione, perché il
GA ottimizza per quella metrica specifica.
\end{domanda}

% ─────────────────────────────────────────────────────────────
\subsection{Domande di approfondimento}
% ─────────────────────────────────────────────────────────────

\begin{domanda}
\textit{``Il GA potrebbe far evolvere anche i pesi della rete?''}\\
Tecnicamente sì, ma sarebbe enormemente inefficiente. Lo spazio dei
pesi è continuo e ad altissima dimensionalità (decine di migliaia di
parametri reali). Il GA è efficiente su spazi discreti e di dimensione
moderata, mentre la backpropagation, che usa l'informazione del
gradiente, è ordini di grandezza più rapida sui pesi. La filosofia del
progetto è usare lo strumento giusto per ciascun livello: GA sui
parametri discreti (architettura), backpropagation sui pesi continui.
\end{domanda}

\begin{domanda}
\textit{``Come potresti migliorare il sistema?''}\\
Diverse strade. Estendere il cromosoma per includere learning rate ed
epochs (oggi fissi). Crossover ``intelligenti'' che scambiano layer simili
in posizione simile. Fitness multi-obiettivo (NSGA-II) per esplorare il
fronte di Pareto fra accuracy e complessità. Surrogate models per
stimare la fitness senza training completo, abbattendo il costo. Estendere
lo spazio di ricerca a topologie non sequenziali (skip connections).
\end{domanda}

\begin{domanda}
\textit{``Il tuo GA è collegato alla ricerca locale del libro Russell-Norvig?''}\\
Sì. Il GA può essere visto come una variante della ricerca locale
stocastica. Hill climbing tiene un singolo stato e va al successore
migliore. Simulated annealing aggiunge tolleranza alle mosse
peggiorative. Local beam search tiene $k$ stati in parallelo. Il GA
estende il beam search aggiungendo il \emph{crossover}: i nuovi
candidati possono combinare informazioni da più stati genitori, non
solo essere successori di uno solo. Aggiunge anche la mutazione, che è
analoga al rumore casuale di simulated annealing.
\end{domanda}

\begin{domanda}
\textit{``Ti aspettavi che il GA trovasse architetture con un solo layer?''}\\
A priori no. Mi aspettavo architetture più profonde, soprattutto su Digits
che ha $64$ feature e $10$ classi. Il GA ha invece dimostrato che per
questi dataset \emph{tabulari} un singolo layer abbastanza ampio è già
sufficiente: l'extra profondità non aggiunge accuracy ma aggiunge
parametri, quindi viene scoraggiata dalla penalità $\lambda$. Questo è
coerente con la letteratura: la profondità è davvero utile quando i dati
hanno struttura gerarchica (immagini, linguaggio), non altrettanto sui
dataset tabulari semplici.
\end{domanda}

\newpage

% =============================================================
\appendix
\section{Glossario dei termini chiave}
% =============================================================

\begin{description}
    \item[Accuracy] Percentuale di predizioni corrette sul set di
        valutazione. Compresa fra $0$ e $1$ (o $0$ e $100\%$).
    \item[Architettura] Struttura della rete neurale: numero di layer,
        neuroni per layer, funzioni di attivazione. È un iperparametro.
        Nel progetto è il cromosoma del GA.
    \item[Backpropagation] Algoritmo per calcolare il gradiente della
        loss rispetto ai pesi della rete, applicando la regola della
        catena layer per layer all'indietro.
    \item[Baseline] Modello di riferimento (qui: rete con architettura
        fissa $[(8,\text{leaky\_relu}),(8,\text{leaky\_relu})]$) usato per
        capire se il GA sta aggiungendo valore.
    \item[Bias (di un neurone)] Termine costante $b$ aggiunto alla somma
        pesata degli input prima dell'attivazione: $f(\bm{w}^\top \bm{x} + b)$.
    \item[Bias ottimistico val/test] Sovrastima sistematica delle
        prestazioni quando si riporta come risultato la stessa metrica
        usata per selezionare il modello.
    \item[Classificazione] Problema in cui $y$ è una categoria discreta.
        Metrica naturale: accuracy. Funzione di output: softmax.
    \item[Cromosoma] Rappresentazione di una soluzione candidata nel GA.
        Nel progetto: lista di tuple \texttt{(n\_neuroni, funzione)}.
    \item[Cross-entropy] Loss function per classificazione. La forma
        categoriale per multiclasse: $-\sum_i y_i \log \hat{y}_i$.
    \item[Crossover] Operatore genetico che combina due genitori per
        generare un figlio.
    \item[Data leakage] Uso involontario di informazione del validation
        o test set durante il training, che gonfia le stime di
        performance.
    \item[Elitismo] Tecnica di copiare il miglior individuo nella
        generazione successiva senza farlo subire crossover/mutazione.
    \item[Esplorazione] Capacità del GA di visitare zone nuove dello
        spazio delle soluzioni. Promossa da mutazione, popolazione
        grande, tournament size piccolo.
    \item[Fitness] Misura di qualità di un individuo nel GA. Nel
        progetto: $\text{accuracy} - \lambda \cdot \log_{10}(N_{\text{params}})$.
    \item[Forward pass] Calcolo dell'output di una rete a partire da un
        input.
    \item[Gene] Elemento atomico di un cromosoma. Nel progetto: una
        tupla \texttt{(n\_neuroni, funzione)}.
    \item[Generazione] Una iterazione completa del ciclo evolutivo del GA.
    \item[Gradient descent] Algoritmo di ottimizzazione che muove i
        parametri nella direzione opposta al gradiente della loss.
    \item[Iperparametri] Parametri di un modello che non vengono appresi
        dai dati ma fissati prima del training (learning rate, $\lambda$,
        ecc.). L'architettura è un iperparametro.
    \item[Learning rate ($\eta$)] Iperparametro che controlla la
        dimensione del passo di gradient descent. Il più critico nel
        progetto.
    \item[Loss] Funzione che misura quanto la predizione si discosta
        dalla verità. La backpropagation cerca di minimizzarla.
    \item[Mutazione] Operatore genetico che introduce variazioni casuali
        in un cromosoma per mantenere diversità.
    \item[NAS (Neural Architecture Search)] Sottocampo dell'AutoML che
        automatizza la progettazione di architetture neurali.
    \item[One-hot encoding] Rappresentazione delle etichette di classe
        come vettori binari con un solo $1$ nella posizione della classe.
    \item[Overfitting] Fenomeno per cui un modello memorizza il training
        invece di generalizzare. Si manifesta con accuracy alta sul
        training ma bassa sul validation/test.
    \item[Parametri (di un modello)] Valori che il modello impara dai
        dati: nel caso di una rete neurale, pesi e bias.
    \item[Penalità di complessità] Termine sottratto alla fitness per
        scoraggiare reti grandi. Nel progetto: $\lambda \cdot
        \log_{10}(N_{\text{params}})$.
    \item[Popolazione] Insieme di tutti gli individui in una generazione.
    \item[ReLU] $f(x) = \max(0, x)$. La funzione di attivazione più
        usata nei layer nascosti.
    \item[Selezione a torneo] Procedura di selezione: si estraggono $k$
        individui casuali e si sceglie il migliore tra loro.
    \item[Softmax] Funzione di output per classificazione multiclasse:
        trasforma un vettore in una distribuzione di probabilità.
    \item[Sfruttamento] Capacità del GA di concentrarsi sulle zone già
        promettenti. Promosso da elitismo, tournament size grande, basso
        mutation rate.
    \item[StandardScaler] Trasformazione che porta ogni feature a media
        $0$ e deviazione standard $1$. Da fittare \emph{solo} sul
        training.
    \item[Stratificazione] Tecnica di split che preserva le proporzioni
        delle classi in ciascun sottoinsieme.
    \item[Test set] Sottoinsieme dei dati toccato \emph{una sola volta}
        alla fine, per stimare onestamente la performance del modello.
    \item[Training set] Sottoinsieme dei dati usato per aggiornare i
        pesi via backpropagation.
    \item[Validation set] Sottoinsieme usato per scegliere fra
        configurazioni alternative (qui: per misurare la fitness del GA).
\end{description}

\end{document}